\documentclass[fontsize=11pt, twoside=false, numbers=autoenddot]{scrbook}
\usepackage{automaten_tafelanschriebe}

\pagestyle{plain}
% \pagestyle{scrheadings}
% % \chead{\headmark}
% \chead{\partmark}
% \renewcommand{\sectionmark}[1]{\markright{\textsl{#1}}}
% \renewcommand{\chaptermark}[1]{\markright{\textsl{#1}}{}}
% \renewcommand{\partmark}[1]{\markright{\textsl{#1}}}
\parindent0pt
\parskip\smallskipamount

\title{Tafelmitschriften zur Vorlesung \glqq Automatentheorie und ihre Anwendungen\grqq\\ im Wintersemester 2018/19}
\author{%
  Prof.\ Dr.\ Thomas Schneider\\[1pt]
  AG Theorie der Künstlichen Intelligenz \\[1pt]
  Fachbereich 3 \\
  \includegraphics[width=.4\linewidth]{logo_ub.jpg} \\[\baselineskip]~%
}
\date{Stand: \today}
\publishers{{\large Dieses Dokument ist noch unvollständig und wird regelmäßig aktualisiert.}}

\begin{document}

\maketitle
\tableofcontents

% ===================================================================
% ===================================================================
% ===================================================================
\part[Endliche Automaten auf endlichen Wörtern]{Endliche Automaten \\ auf endlichen Wörtern}

%\dots\ folgt \dots

% ===================================================================
% ===================================================================
% ===================================================================
\part[Endliche Automaten auf endlichen Bäumen]{Endliche Automaten \\ auf endlichen Bäumen}

%\dots\ folgt \dots

% ===================================================================
% ===================================================================
% ===================================================================
\part[Endliche Automaten auf unendlichen Wörtern]{Endliche Automaten \\ auf unendlichen Wörtern}

% ===================================================================
\section*{T3.1~ Produktkonstruktion für NEAs ist für NBAs nicht korrekt}

Sei $\Sigma = \{a,b\}$. Wir betrachten folgende \textsfbf{NEAs} $\Amc_1,\Amc_2$.
%
\begin{center}
  \begin{tikzpicture}[%
    node distance=20mm,>=Latex,
    initial text="", initial where=below left,
    every state/.style={draw=black,thin,fill=black!5,inner sep=1mm,minimum size=6mm},
    accepting/.style={double distance=1.5pt, double=white},
    every edge/.style={draw=black,thin}
  ]
    \node[state,initial]   (q0) {$q_0$};
    \node[state,accepting] (q1) [right of=q0] {$q_1$};
    \node[above left=-2mm and 6mm of q0] {$\Amc_1$};
    
    \path[->] (q0) edge [bend left=15] node [above] {$a$} (q1)
              (q1) edge [bend left=15] node [below] {$a,b$} (q0);
  \end{tikzpicture}
  \qquad
  \begin{tikzpicture}[%
    node distance=20mm,>=Latex,
    initial text="", initial where=below left,
    every state/.style={draw=black,thin,fill=black!5,inner sep=1mm,minimum size=6mm},
    accepting/.style={double distance=1.5pt, double=white},
    every edge/.style={draw=black,thin}
  ]
    \node[state,initial,accepting] (q2) {$q_2$};
    \node[state]                   (q3) [right of=q2] {$q_3$};
    \node[above left=-2mm and 6mm of q2] {$\Amc_2$};
    
    \path[->] (q2) edge [bend left=15] node [above] {$a,b$} (q3)
              (q3) edge [bend left=15] node [below] {$b$}   (q2);
  \end{tikzpicture}
\end{center}
%
Dann gilt:
%
\begin{align*}
  L(\Amc_1) & = \{a_0\cdots a_{n-1} \mid n \text{~ist ungerade und~} a_0=a_2=\dots=a\} \\
  L(\Amc_2) & = \{a_0\cdots a_{n-1} \mid n \text{~ist gerade und~} a_1=a_3=\dots=b\} \\[6pt]
  L(\Amc_1) \cap L(\Amc_2) & = \emptyset
\end{align*}
%
Der Produktautomat \Amc ist folgender:
%
\begin{center}
  \begin{tikzpicture}[%
    node distance=40mm,>=Latex,
    initial text="", initial where=below left,
    every state/.style={ellipse,draw=black,thin,fill=black!5,inner sep=1mm,minimum size=6mm},
    accepting/.style={double distance=1.5pt, double=white},
    every edge/.style={draw=black,thin}
  ]
    \node[state,initial]   (q0q2) {$(q_0,q_2)$};
    \node[state,accepting] (q1q2) [right of=q0q2]      {$(q_1,q_2)$};
    \node[state]           (q0q3) [below=10mm of q0q2] {$(q_0,q_3)$};
    \node[state]           (q1q3) [right of=q0q3]      {$(q_1,q_3)$};
    \node[above left=-2mm and 10mm of q0q2] {$\Amc$};
    
    \path[->] (q0q2) edge [bend left=10] node [above,near start] {$a$} (q1q3)
              (q1q3) edge [bend left=10] node [below,near end] {$b$} (q0q2)
              (q1q2) edge [bend left=10] node [below,pos=.09] {$\,a,b$} (q0q3);
  \end{tikzpicture}
\end{center}
%
Diese Konstruktion ist korrekt für NEAs; in diesem Beispiel ist $L(\Amc) = \emptyset$,
da der einzige akzeptierende Zustand $(q_1,q_2)$ unerreichbar ist.

\par\medskip\noindent
Betrachten wir nun dieselben Automaten $\Amc_1,\Amc_2$ als \textsfbf{NBAs.}
Dann gilt:
%
\begin{align*}
  L_\omega(\Amc_1) & = \{\alpha \mid n \text{~ist ungerade und~} \alpha_0=\alpha_2=\dots=a\} \\
  L_\omega(\Amc_2) & = \{\alpha \mid n \text{~ist gerade und~} \alpha_1=\alpha_3=\dots=b\},
  \intertext{also ist jetzt}
  L_\omega(\Amc_1) \cap L_\omega(\Amc_2) & = \{(ab)^\omega\},
  \intertext{aber nach wie vor ist}
  L_\omega(\Amc) = \emptyset,
\end{align*}
%
also ist die Konstruktion für NBAs nicht korrekt!
Der Grund dafür ist, dass die erfolgreichen Runs von $\Amc_1$ bzw.\ $\Amc_2$
die akzeptierenden Zustände asynchron erreichen,
nämlich nach $1,3,5,\dots$ Schritten $(\Amc_1)$ bzw.\
$0,2,4,\dots$ Schritten $(\Amc_2)$. Dadurch erreicht der entsprechende Run
von $\Amc$ niemals einen (kombinierten) akzeptierenden Zustand.

\goodbreak
% ===================================================================
\section*{T3.2~ Produktkonstruktion für NBAs: Beispiel und Korrektheit}

\paragraph*{Beispiel.}
Wendet man die Produktkonstruktion (Folie~30) auf die obigen NBAs
$\Amc_1,\Amc_2$ an, so erhält man folgenden NBA $\Amc$
(im Bild sind nur die erreichbaren Zustände wiedergegeben):
%
\begin{center}
  \begin{tikzpicture}[%
    node distance=40mm,>=Latex,
    initial text="", initial where=below left,
    every state/.style={ellipse,draw=black,thin,fill=black!5,inner sep=1mm,minimum size=6mm},
    accepting/.style={double distance=1.5pt, double=white},
    every edge/.style={draw=black,thin}
  ]
    \node[state,initial]   (q0q21)                       {$(q_0,q_2,1)$};
    \node[state]           (q1q31) [right of=q0q21]      {$(q_1,q_3,1)$};
    \node[state,accepting] (q0q22) [below=10mm of q1q31] {$(q_0,q_2,2)$};
    \node[above left=-2mm and 12mm of q0q21] {$\Amc$};
    
    \path[->] (q0q21) edge node [above] {$a$} (q1q31)
              (q1q31) edge [bend left=15] node [right] {$b$} (q0q22)
              (q0q22) edge [bend left=15] node [left]  {$a$} (q1q31);
  \end{tikzpicture}
\end{center}
%
Tatsächlich ist nun $L_\omega(\Amc) = \{(ab)^\omega\}$.

\paragraph*{Korrektheitsbeweis.}
Zu zeigen ist:
\[
  L_\omega(\Amc) = L_\omega(\Amc_1) \cap L_\omega(\Amc_2)
\]
%\begin{beweis}
  \begin{description}
    \item[{\boldmath"`$\subseteq$"'}]
      Sei $\alpha \in L_\omega(\Amc)$.
      Dann gibt es einen erfolgreichen Run $r = q_0q_1q_2\cdots$ von \Amc auf $\alpha$
      mit $q_0 \in I$ und $\textsf{Inf}(r) \cap F \neq \emptyset$.
      Nach Konstruktion von \Amc muss jedes $q_i$ die Form
      \[
        q_i = (s_i,t_0,n_i)
      \]
      haben mit $s_i \in Q_1$, $t_i \in Q_2$ und $n_i \in \{1,2\}$ für alle $i \geq 0$.
      
      Wir betrachten die Folge $s = s_0s_1s_2\cdots$.
      Diese ist ein Run von $\Amc_1$ auf $\alpha$:
      da $r$ ein Run ist, folgt mit der Definition von $I$ bzw.\ $\Delta$, dass
      $s_0 \in I_1$ und $(s_i,\alpha_i,s_{i+1}) \in \Delta_1$ für alle $i \geq 0$.
      Außerdem ist $s$ erfolgreich, denn wegen $\textsf{Inf}(r) \cap F \neq \emptyset$
      (und der Definition von $F$)
      enthält $r$ unendlich viele Zustände der Form
      \[
        \tag{$*$}
        q_i = (s_i,t_i,2) \qquad \text{mit~} t_i \in F_2;
      \]
      also enthält $r$ auch unendlich viele Zustände der Form
      \[
        \tag{$**$}
        q_j = (s_j,t_j,1) \qquad \text{mit~} s_j \in F_1,
      \]
      weil nach jedem $q_i$ der Form $(*)$ in $(s_{i+1},t_{i+1},1)$ gewechselt wird
      und erst dann wieder ins nächste $q_i$ der Form $(*)$ gegangen werden kann,
      wenn ein $q_j$ der Form $(**)$ gefunden wurde.
      Also ist $\textsf{Inf}(s) \cap F_1 \neq \emptyset$ und damit $s$ erfolgreich.
      
      Analog argumentiert man, dass $t = t_0t_1t_2\cdots$
      ein erfolgreicher Run von $\Amc_2$ auf $\alpha$ ist.
      Folglich ist $\alpha \in L_\omega(\Amc_1) \cap L_\omega(\Amc_2)$.
    \item[{\boldmath"`$\supseteq$"'}]
      Sei $\alpha \in L_\omega(\Amc_1) \cap L_\omega(\Amc_2)$.
      Dann gibt es erfolgreiche Runs
      %
      \begin{align*}
        s & = s_0s_1s_2\dots \quad \text{von $\Amc_1$ auf $\alpha$\quad und} \\
        t & = t_0t_1t_2\dots \quad \text{von $\Amc_2$ auf $\alpha$.}
      \end{align*}
      %
      Wir betrachten die Folge
      \[
        r = (s_0,t_0,n_0)~(s_1,t_1,n_1)~(s_2,t_2,n_2)~\cdots,
      \]
      wobei die $n_i$ induktiv wie folgt definiert sind:
      %
      \begin{align*}
        n_0 & = 1 \\
        n_i & = \begin{cases}
                  1 & \text{falls~} n_{i-1} = 1 \text{~und~} s_{i-1} \notin F_1 \\
                    & \text{oder~}  n_{i-1} = 2 \text{~und~} t_{i-1} \in F_2 \\
                  2 & \text{sonst}
                \end{cases}
      \end{align*}
      %
      Man zeigt leicht unter Zuhilfenahme der Konstruktion von $I,F,\Delta$,
      dass $r$ ein erfolgreicher Run von \Amc auf $\alpha$ ist.
      Folglich ist $\alpha \in L_\omega(\Amc)$.
      \qedhere
  \end{description}
%\end{beweis}

% ===================================================================
\section*{T3.3~ Büchi-Erkennbarkeit von {\boldmath $W^\omega$ für reguläre Sprachen $W$}}

\textsfbf{Noch zu zeigen:}~ $L_\omega(\Amc_2) = L(\Amc_1)^\omega$
%
\begin{description}
  \item[{\boldmath"`$\subseteq$"'}]
    Sei $\alpha \in L_\omega(\Amc_2)$.
    Dann gibt es einen erfolgreichen Run $r = q_0q_1q_2\cdots$ von $\Amc_2$
    auf $\alpha$, d.\,h.\ $q_0 = q_I$ (der einzige Anfangszustand von $\Amc_2$),
    und $q_I$ kommt unendlich oft in $r$ vor (weil es auch der einzige akzeptierende Zustand von $\Amc_2$ ist).
    
    Seien $q_{i_0},q_{i_1},q_{i_2},\dots$ alle Vorkommen von $q_I$ in $r$.
    Für jedes $j \geq 0$ betrachten wir die Folge
    \[
        r_j := \underset{\substack{||\\q_I}}{q_{i_j}}
               q_{i_j+1}\cdots q_{i_{j+1}-1}
               \underset{\substack{||\\q_I}}{q_{i_{j+1}}}.
    \]
    Nach Konstruktion von $\Delta_2$ und wegen der Annahmen über $\Amc_1$
    gibt es ein $q_f \in F$, so dass
    \[
        \underset{\substack{||\\q_I}}{q_{i_j}}
        q_{i_j+1}\cdots q_{i_{j+1}-1}
        \underset{\substack{\turnbox{270}{$\in$}\,\,\,\\[7pt]F}}{q_f}
    \]
    ein erfolgreicher Run von $\Amc_1$ auf
%    $w_j := \alpha_{i_j}\alpha_{i_j+1}\cdots\alpha_{i_{j+1}-1}$
    $w_j := \alpha[i_j,i_{j+1}-1]$
    ist.
    Folglich gehört für alle $j$ das Wort $w_j$ zur Sprache $L(\Amc_1)$,
    und damit gilt $\alpha \in L(\Amc_1)^\omega$.
  \item[{\boldmath"`$\supseteq$"'}]
    Sei $\alpha \in L(\Amc_1)^\omega$.
    Dann ist $\alpha = w_0w_1w_2\cdots$ mit $w_i \in L(\Amc_1)$ für alle $i \geq 0$.
    Wir nehmen o.\,B.\,d.\,A.\ an, dass $\varepsilon \notin L(\Amc_1)$ ist,
    also $|w_j| > 0$ für alle $i$.
    Also gibt es für jedes $j \geq 0$ einen erfolgreichen Run
    \[
      r_j := q_{j,0}q_{j,1}\cdots q_{j,|w_j|}
    \]
    von $\Amc_1$ auf $w_j$.
    Nach Konstruktion von $\Delta_2$ ist dann
    \[
      r ~:=~ q_{0,0}q_{0,1}\cdots q_{0,|w_0|-1}~~
             q_{1,0}q_{1,1}\cdots q_{1,|w_1|-1}~~
             \cdots
    \]
    ein erfolgreicher Run von $\Amc_2$ auf $\alpha$.
    Folglich ist $\alpha \in L_\omega(\Amc_2)$.
    \qedhere
\end{description}

% ===================================================================
\section*{T3.4~ Beweis Charakterisierung NBA-erkennbarer Sprachen}

\textsfbf{Satz~3.9.}~
Eine Sprache $L \subseteq \Sigma^\omega$ ist Büchi-erkennbar
genau dann,
wenn es reguläre Sprachen $V_1,W_1,\dots,V_n,W_n$ gibt mit $n \geqslant 1$ und
\[
  L = V_1W_1^\omega \cup \dots \cup V_nW_n^\omega.
\]

\par\medskip
\textsfbf{Beweis.}~
Sei $L \subseteq \Sigma^\omega$ Büchi-erkennbar,
also $L = L_\omega(\Amc)$ für einen NBA $\Amc = (Q,\Sigma,\Delta,I,F)$.
Wir nutzen folgende Beobachtung: für jedes $\omega$-Wort $\alpha \in L_\omega(\Amc)$,
auf dem $\Amc$ einen erfolgreichen Run $r = q_0q_1q_2\cdots$
mit $q_f \in \textsf{Inf}(r) \cap F$ hat, gibt es
%
\begin{itemize}
  \item
    ein endliches Präfix von $\alpha$, das $q_I$ nach $q_f$ überführt und
  \item
    unendlich viele nicht-leere Infixe, die $q_f$ nach $q_f$ überführen.
\end{itemize}
%
Diese beiden Sorten von Infixen können wir durch reguläre Sprachen beschreiben.
Dazu verwenden wir folgende Notation: Für zwei beliebige Zustände $q_1,q_2 \in Q$
sei $\Amc_{q_1,q_2} = (Q,\Sigma,\Delta,\{q_1\},\{q_2\})$
und $W_{q_1,q_2} = L(\Amc_{q_1,q_2})$.
Nach Definition sind die $W_{q_1,q_2}$ regulär.
Wegen der Akzeptanzbedingung von Büchi-Automaten und unserer Beobachtung gilt nun:
\[
  L_\omega = \bigcup_{\substack{q_i \in I\\q_f \in F}} W_{q_i,q_f} W_{q_f,q_f}^\omega
\]
\qedhere

% ===================================================================
\section*{T3.5~ Beispiele für {\boldmath $\Vec W$}}

\begin{align*}
  W_1       & = \{a^nb^m \mid n,m \geq 0\}                                           \\
  \Vec{W_1} & = \{a^nb^\omega \mid n \geq 0\} \cup \{a^\omega\}                      \\[8pt]
  W_2       & = \{a^nb^n \mid n \geq 0\}                                             \\
  \Vec{W_2} & = \emptyset                                                            \\[8pt]
  W_3       & = \{a,b\}^*                                                            \\
  \Vec{W_3} & = \{a,b\}^\omega                                                       \\[8pt]
  W_4       & = \{w \in \{a,b\}^* \mid \#_a(w) \text{~ist gerade}\}                  \\
  \Vec{W_4} & = \{\alpha \in \{a,b\}^\omega \mid \#_a(\alpha) = \infty \text{~oder~}
                 \alpha = wb^\omega \text{~mit~} \#_a(w) \text{~gerade}\}            \\[8pt]
  W_5       & = \{w \in \{a,b\}^* \mid \#_a(w) = \#_b(w)\}                           \\
  \Vec{W_5} & = \Big\{\alpha \in \{a,b\}^\omega ~\Big|~ \#\{i \mid \#_a(\alpha[0,i]) = \#_b(\alpha[0,i])\} = \infty\Big\}            \\[8pt]
\end{align*}

% ===================================================================
\section*{T3.6~ Beweis Charakterisierung DBA-erkennbarer Sprachen}

\textsfbf{Satz~3.11.}~
Eine $\omega$-Sprache $L \subseteq \Sigma^\omega$
ist DBA-erkennbar genau dann,
wenn es eine reguläre Sprache $W \subseteq \Sigma^*$ gibt
mit $L = \Vec W$.

\par\medskip
\textsfbf{Beweis.}~
Es genügt zu zeigen, dass
für jeden \textsfbf{D}EA/\textsfbf{D}BA $\Amc=(Q,\Sigma,\Delta,\{q_I\},F)$ gilt:
\[
  L_\omega(\Amc) = \Vec{L(\Amc)}
\]
\begin{description}
  \item[{\boldmath"`$\subseteq$"'}]
    (Diese Richtung funktioniert sogar, wenn $\Amc$ ein \textsfbf{N}EA/\textsfbf{N}BA ist.)
    
    Sei $\alpha \in L_\omega(\Amc)$.
    Dann gibt es einen erfolgreichen Run $r = q_0q_1q_2\cdots$
    von $\Amc$ auf $\alpha$.
    Seien $i_0,i_1,i_2,\ldots \in \mathbb{N}$ die Positionen mit $q_{i_j} \in F$.
    Dann ist für jedes $j \geq 0$ das Präfix $q_0\cdots q_{i_j}$ von $r$ ein
    erfolgreicher Run des \textsfbf{NEAs} $\Amc$ auf $\alpha_0\cdots\alpha_{i_j-1}$.
    Damit gibt es unendlich viele Präfixe von $\alpha$, die in $L(\Amc)$ sind,
    und damit ist $\alpha \in \Vec{L(\Amc)}$.
  \item[{\boldmath"`$\supseteq$"'}]
    Sei $\alpha \in \Vec{L(\Amc)}$.
    Dann hat $\alpha$ unendlich viele Präfixe in $L(\Amc)$.
    Also muss der \textsfbf{eindeutig bestimmte} Run
    $r$ von $\Amc$ (der Automat ist deterministisch!)
    einen akzeptierenden Zustand unendlich oft erreichen.
    Damit ist $\alpha \in L_\omega(\Amc)$.
    \qedhere
\end{description}

% ===================================================================
\section*{T3.7~ Beispiele für Muller-Automaten}

\begin{itemize}
  \item 
    Betrachte den folgenden Muller-Automaten $\Amc_1 = (Q_1,\Sigma,\Delta_1,I_1,\Fmc_1)$ über dem Alphabet $\Sigma = \{a,b\}$.
    %
    \begin{center}
      \begin{tikzpicture}[%
        node distance=25mm,>=Latex,
        initial text="", initial where=below left,
        every state/.style={draw=black,thin,fill=black!5,inner sep=1mm,minimum size=6mm},
        accepting/.style={double distance=1.5pt, double=white},
        every edge/.style={draw=black,thin}
      ]
        \node[state,initial]   (q0) {$q_0$};
        \node[state]           (q1) [right of=q0] {$q_1$};
        \node[above left=-2mm and 6mm of q0] {$\Amc_1$};
        
        \path[->] (q0) edge node [above] {$b$} (q1)
                  (q0) edge [loop above] node [right] {$~a,b$} ()
                  (q1) edge [loop above] node [right] {$~b$} ();
      \end{tikzpicture}
    \end{center}
    %
    \begin{itemize}
      \item[\bspnr{(M1)}]
        Wenn $\Fmc_1 = \{\{q_0\}\}$, dann $L_\omega(\Amc_1) = \Sigma^\omega$.
      \item[\bspnr{(M2)}]
        Wenn $\Fmc_1 = \{\{q_1\}\}$, dann $L_\omega(\Amc_1) = \{\alpha \in \Sigma^\omega \mid \#_a(\alpha) < \infty\}$.
      \item[\bspnr{(M3)}]
        Wenn $\Fmc_1 = \{\{q_0,q_1\}\}$, dann $L_\omega(\Amc_1) = \emptyset$ (weil Wechsel von $q_1$ zu $q_0$ nicht möglich).
      \item[\bspnr{(M4)}]
        Wenn $\Fmc_1 = \{\{q_0\},\{q_1\}\}$, dann $L_\omega(\Amc_1) = \Sigma^\omega$ (Vereinigung der Fälle \bspnr{M1},\,\bspnr{M2}).
    \end{itemize}
    \par\medskip    
  \item 
    Betrachte nun den folgenden Muller-Automaten $\Amc_2 = (Q_2,\Sigma,\Delta_2,I_2,\Fmc_2)$ über demselben Alphabet.
    %
    \begin{center}
      \begin{tikzpicture}[%
        node distance=25mm,>=Latex,
        initial text="", initial where=below left,
        every state/.style={draw=black,thin,fill=black!5,inner sep=1mm,minimum size=6mm},
        accepting/.style={double distance=1.5pt, double=white},
        every edge/.style={draw=black,thin}
      ]
        \node[state,initial] (q0) {$q_0$};
        \node[state]         (q1) [right of=q2] {$q_1$};
        \node[above left=-2mm and 6mm of q2] {$\Amc_2$};
        
        \path[->] (q0) edge [bend left=15] node [above] {$b$} (q1)
                  (q1) edge [bend left=15] node [below] {$b$} (q0)
                  (q0) edge [loop above]   node [right] {$~a$} ();
      \end{tikzpicture}
    \end{center}
    %
    \begin{itemize}
      \item[\bspnr{(M5)}]
        Wenn $\Fmc_2 = \{\{q_0\}\}$, dann $L_\omega(\Amc_2) = L((a+bb)^*a^\omega)$ (Menge aller Wörter mit endlich vielen $b$'s, in denen zwischen je zwei $a$'s und vor dem ersten $a$ eine gerade Anzahl von $b$'s steht).
      \item[\bspnr{(M6)}]
        Wenn $\Fmc_2 = \{\{q_1\}\}$, dann $L_\omega(\Amc_2) = \emptyset$ (denn wenn $q_1$ unendlich oft besucht wird, dann auch $q_0$).
      \item[\bspnr{(M7)}]
        Wenn $\Fmc_2 = \{\{q_0,q_1\}\}$, dann $L_\omega(\Amc_2) = L((a^*bb)^\omega)$ (Menge aller Wörter wie in \bspnr{M5}, aber mit \emph{un}endlich vielen $b$'s).
    \end{itemize}
    \par\medskip    
  \item 
    Für \emph{alle} Muller-Automaten $\Amc = (Q,\Sigma,\Delta,I,\Fmc)$ gilt:
    wenn $\Fmc = \emptyset$ oder $\Fmc = \{\emptyset\}$,
    dann $L_\omega(\Amc) = \emptyset$.
\end{itemize}

% ===================================================================
\section*{T3.8~ Beispiele für Rabin-Automaten}

\begin{itemize}
  \item 
    Betrachte den folgenden Rabin-Automaten $\Amc = (Q,\Sigma,\Delta,I,\Pmc)$ über dem Alphabet $\Sigma = \{a,b\}$.
    %
    \begin{center}
      \begin{tikzpicture}[%
        node distance=25mm,>=Latex,
        initial text="", initial where=below left,
        every state/.style={draw=black,thin,fill=black!5,inner sep=1mm,minimum size=6mm},
        accepting/.style={double distance=1.5pt, double=white},
        every edge/.style={draw=black,thin}
      ]
        \node[state,initial] (q0) {$q_0$};
        \node[state]         (q1) [right of=q0] {$q_1$};
        \node[above left=-2mm and 6mm of q0] {$\Amc_2$};
        
        \path[->] (q0) edge [bend left=15] node [above] {$b$} (q1)
                  (q1) edge [bend left=15] node [below] {$b$} (q0)
                  (q0) edge [loop above]   node [right] {$~a$} ();
      \end{tikzpicture}
    \end{center}
    %
    \begin{itemize}
      \item[\bspnr{(R1)}]
        Wenn $\Pmc = \{(\{q_0\},\{q_1\})\}$, dann $L_\omega(\Amc) = \emptyset$ (Begründung wie bei \bspnr{M6}).
      \item[\bspnr{(R2)}]
        Wenn $\Pmc = \{(\{q_1\},\{q_0\})\}$, dann $L_\omega(\Amc) = L((a+bb)^*a^\omega)$ (dasselbe Akzeptanzverhalten wie in \bspnr{M5}).
      \item[\bspnr{(R3)}]
        Wenn $\Pmc = \{(\emptyset,\{q_1\})\}$, dann $L_\omega(\Amc) = L((a^*bb)^\omega)$ (dasselbe Akzeptanzverhalten wie NBA mit $\Fmc = \{q_1\}$).
      \item[\bspnr{(R4)}]
        Wenn $\Pmc = \{\{(S,\emptyset)\}\}$ für beliebiges $S \subseteq Q$, dann $L_\omega(\Amc) = \emptyset$ (folgt direkt aus Definition "`erfolgreich"' für NRAs).
    \end{itemize}
    \par\medskip    
  \item 
    Der Fall mehrerer Paare in der Akzeptanzkomponente $\Pmc$ braucht nicht gesondert illustriert zu werden,
    denn für \emph{alle} NRAs $\Amc = (Q,\Sigma,\Delta,I,\Pmc)$ mit $\Pmc = \bigcup_{i \leq n} \Pmc_i$ gilt:
    $L_\omega(\Amc) = \bigcup_{i \leq n} L_\omega(Q,\Sigma,\Delta,I,\Pmc_i)$.
\end{itemize}

% ===================================================================
\section*{T3.9~ Beispiele für Streett-Automaten}

\begin{itemize}
  \item 
    Betrachte den folgenden Streett-Automaten $\Amc = (Q,\Sigma,\Delta,I,\Pmc)$ über dem Alphabet $\Sigma = \{a,b\}$.
    %
    \begin{center}
      \begin{tikzpicture}[%
        node distance=25mm,>=Latex,
        initial text="", initial where=below left,
        every state/.style={draw=black,thin,fill=black!5,inner sep=1mm,minimum size=6mm},
        accepting/.style={double distance=1.5pt, double=white},
        every edge/.style={draw=black,thin}
      ]
        \node[state,initial] (q0) {$q_0$};
        \node[state]         (q1) [right of=q0] {$q_1$};
        \node[above left=-2mm and 6mm of q0] {$\Amc_2$};
        
        \path[->] (q0) edge [bend left=15] node [above] {$b$} (q1)
                  (q1) edge [bend left=15] node [below] {$b$} (q0)
                  (q0) edge [loop above]   node [right] {$~a$} ();
      \end{tikzpicture}
    \end{center}
    %
    \begin{itemize}
      \item[\bspnr{(S1)}]
        Wenn $\Pmc = \{(\{q_0\},\{q_1\})\}$, dann $L_\omega(\Amc) = (a+bb)^\omega$ (denn jeder Run ist erfolgreich).
      \item[\bspnr{(S2)}]
        Wenn $\Pmc = \{(\{q_0\},\emptyset)\}$, dann $L_\omega(\Amc) = (a+bb)^\omega$ (denn jeder Run ist erfolgreich).
      \item[\bspnr{(S3)}]
        Wenn $\Pmc = \{(\{q_1\},\{q_0\})\}$, dann $L_\omega(\Amc) = (a*bb)^\omega$ (dasselbe Akzeptanzverhalten wie \bspnr{M7}).
      \item[\bspnr{(S4)}]
        Wenn $\Pmc = \{(\emptyset,\{q_1\})\}$, dann $L_\omega(\Amc) = (a+bb)*a^\omega$ (denn die Akzeptanzbedingung besagt: $q_1$ darf \emph{nicht} $\infty$ oft vorkommen, also muss $q_0$ $\infty$ oft vorkommen $\leadsto$ wie \bspnr{M5}).
      \item[\bspnr{(S5)}]
        Wenn $\Pmc = \{(\emptyset,\{q_0\})\}$, dann $L_\omega(\Amc) = \emptyset$ (denn $q_0$ kommt in jedem Run $\infty$ oft vor).
      \item[\bspnr{(S6)}]
        Wenn $\Pmc = \{(\emptyset,\{q_0,q_1\})\}$, dann $L_\omega(\Amc) = \emptyset$ (wie \bspnr{S5}).
    \end{itemize}
    \par\medskip    
  \item 
    Der Fall mehrerer Paare in der Akzeptanzkomponente $\Pmc$ ist analog zu NRAs,
    aber mit einem entscheidenden Unterschied:
    für alle NSAs $\Amc = (Q,\Sigma,\Delta,I,\Pmc)$ mit $\Pmc = \bigcup_{i \leq n} \Pmc_i$ gilt:
    $L_\omega(\Amc) = \bigcap_{i \leq n} L_\omega(Q,\Sigma,\Delta,I,\Pmc_i)$.
\end{itemize}

% ===================================================================
\section*{T3.10~ Überblick Beweis der Gleichmächtigkeit}

Das folgende Bild illustriert, wie die Aussagen der Lemmata 3.17--3.19
zur Äquivalenz der vier Automatenmodelle (Satz 3.16) beitragen.
Ein Pfeil vom Knoten N$x$A zum Knoten N$y$A bedeutet dabei:
"`jede N$x$A-erkennbare Sprache ist N$y$A-erkennbar"';
die Beschriftung der Pfeile gibt die Nummer des Lemmas an.
%
\begin{center}
  \begin{tikzpicture}[%
    node distance=20mm,>=Latex,
    initial text="", initial where=below left,
    every state/.style={rectangle,rounded corners,draw=black,thin,fill=black!5,inner sep=1mm,minimum size=6mm},
    every edge/.style={draw=black,thin}
  ]
    \node[state] (NRA)                   {NRA};
    \node[state,above left of=NRA] (NBA) {NBA};
    \node[state,below left of=NRA] (NSA) {NSA};
    \node[state,right=20mm of NRA]      (NMA) {NMA};
    
    \path[->] (NBA) edge                node [above]            {3.17} (NMA)
              (NRA) edge                node [above,near start] {3.17} (NMA)
              (NSA) edge                node [below]            {~3.17} (NMA)
              (NBA) edge[bend left=10]  node [left,near end]    {3.18} (NRA)
              (NBA) edge[bend right=10] node [left]             {3.18} (NSA)
              (NMA) edge[bend right=30] node [right,near start] {~3.19} (NBA)
              ;
  \end{tikzpicture}
\end{center}

% ===================================================================
\section*{T3.11~ Beweis Korrektheit "`von Muller- zu Büchi-Automaten"'}

Am Ende des Beweises von Lemma~3.19 ist zu zeigen:~ $L_\omega(\Amc') = L_\omega(\Amc)$.
%
\begin{description}
  \item[{\boldmath"`$\supseteq$"'}]
    Sei $\alpha \in L_\omega(\Amc)$ und $r = q_0q_1q_2\cdots$ ein
    erfolgreicher Run von \Amc (NMA!) auf $\alpha$, also $q_0 \in I$
    und $\textsf{Inf}(r) = F$.
    Dann gibt es eine Position $i_0 \geq 1$, ab der nur noch
    akzeptierende Zustände auftreten, d.\,h.\ $q_i \in F$ für alle $i \geq i_0$.
    Daraus konstruieren wir wie folgt induktiv eine Folge $s=s_0s_1s_2\cdots$ von Zuständen
    von $\Amc'$:
    %
    \begin{itemize}
      \item
        Für alle $i < i_0$ setze $s_i = q_i$.
      \item
        $s_{i_0} = \auf q_{i_0},\{q_{i_0}\}\zu$
      \item
        Für alle $i > i_0$ setze
        \[
          s_i = \begin{cases}
                  \auf q_i,S \cup \{q_i\}\zu & \text{falls~} s_{i-1} = \auf q_{i-1},S\zu \text{~und~} S \neq F \\
                  \auf q_i, \{q_i\}\zu & \text{falls~} s_{i-1} = \auf q_{i-1},F\zu 
                \end{cases}
        \]
    \end{itemize}
    %
    Diese Konstruktion von $s$ stellt sicher:
    %
    \begin{itemize}
      \item 
        $s$ ist Run von $\Amc'$ auf $\alpha$ (das ist leicht schrittweise anhand der Konstruktion von $\Delta'$
        und von $s$ nachvollziehbar).
      \item 
        $s_0 \in I$ (nach Definition $I'$).
      \item 
        $s$ ist erfolgreich, denn wegen $\textsf{Inf}(r) = F$
        gibt es $\infty$ viele Vorkommen von Zuständen der Form $\auf q_f,F\zu$ in $S$.
    \end{itemize}
    %
    Also ist $\alpha \in L_\omega(\Amc')$.
  \item[{\boldmath"`$\subseteq$"'}]
    Sei $\alpha \in L_\omega(\Amc')$ und $s = s_0s_1s_2\cdots$ ein
    erfolgreicher Run von $\Amc'$ (NBA!) auf $\alpha$, also $s_0 \in I'$
    und ein $\auf q_f,F\zu$ kommt $\infty$ oft in $s$ vor.
    Konstruiere daraus eine Folge $r = q_0q_1q_2\cdots$ von Zuständen aus $Q$ wie folgt:
    %
    \goodbreak
    \begin{itemize}
      \item
        Wenn $s_i \in Q$, dann $q_i = s_i$.
      \item
        Wenn $s_i = \auf q_f,S\zu$, dann $q_i = q_f$.
    \end{itemize}
    %
    Diese Konstruktion stellt sicher:
    %
    \begin{itemize}
      \item
        $r$ ist Run von $\Amc$ auf $\alpha$ (vgl.\ Konstruktion von $\Delta'$).
      \item
        $r_0 \in I$.
      \item
        $r$ enthält nur endlich viele Zustände außerhalb $F$ (weil in "`Phase 2"' nur noch Zustände aus $F$ vorkommen).
      \item 
        Jeder Zustand aus $F$ kommt in $r$ unendlich oft vor (weil $\auf q_f,F\zu$ unendlich oft in $s$ vorkommt).
    \end{itemize}
    %
    Damit ist $r$ ein erfolgreicher Run des NMA $\Amc$ auf $\alpha$,
    also $\alpha \in L_\omega(\Amc)$.
    \qedhere
\end{description}

% ===================================================================
\section*{T3.12~ Determinisierungsversuch mittels Potenzmengenkonstruktion}

Wir betrachten folgenden NBA $\Amc$ über dem Alphabet $\Sigma = \{a,b\}$.
%
\begin{center}
  \begin{tikzpicture}[%
    node distance=25mm,>=Latex,
    initial text="", initial where=below left,
    every state/.style={draw=black,thin,fill=black!5,inner sep=1mm,minimum size=6mm},
    accepting/.style={double distance=1.5pt, double=white},
    every edge/.style={draw=black,thin}
  ]
    \node[state,initial,accepting] (q0) {$q_0$};
    \node[state]                   (q1) [right of=q0] {$q_1$};
    \node[above left=-2mm and 8mm of q0] {$\Amc$};
    
    \path[->] (q0) edge [bend left=15] node [above] {$b$} (q1)
              (q1) edge [bend left=15] node [below] {$b$} (q0)
              (q1) edge [loop right]   node [right] {$a,b$} ();
  \end{tikzpicture}
\end{center}
%
Die erkannte Sprache ist
$L_\omega(\Amc) = L((b\Sigma^*b)^\omega)
= \{\alpha \in \Sigma^\omega \mid a_0 = b \text{~und~} \#_{bb}(\alpha) = \infty\}$.
Mittels Potenzmengenkonstruktion erhalten wir folgenden DBA $\Amc^d$ (der Papierkorbzustand ist weggelassen).
%
\begin{center}
  \begin{tikzpicture}[%
    node distance=40mm,>=Latex,
    initial text="", initial where=below left,
    every state/.style={ellipse,draw=black,thin,fill=black!5,inner sep=1mm,minimum size=6mm},
    accepting/.style={double distance=1.5pt, double=white},
    every edge/.style={draw=black,thin}
  ]
    \node[state,initial,accepting] (q0)                      {$\{q_0\}$};
    \node[state]                   (q1)   [right of=q0]      {$\{q_1\}$};
    \node[state,accepting]         (q0q1) [below=10mm of q1] {$\{q_0,q_1\}$};
    \node[above left=-2mm and 8mm of q0q21] {$\Amc^d$};
    
    \path[->] (q0)   edge                 node [above] {$b$} (q1)
              (q1)   edge [bend right=15] node [left]  {$b$} (q0q1)
              (q0q1) edge [bend right=15] node [right] {$a$} (q1)
              (q1)   edge [loop right]    node [right] {$a$} ()
              (q0q1) edge [loop right]    node [right] {$b$} ();
  \end{tikzpicture}
\end{center}
%
Nun ist aber $(ba)^\omega \in L_\omega(\Amc^d) \setminus L_\omega(\Amc)$.
Der DBA $\Amc^d$ hat also auf dem Wort $(ba)^\omega$ einen \emph{Bad Run} $r$,
der keinem erfolgreichen Run von $\Amc$ auf $(ba)^\omega$ entspricht.
Der Grund dafür ist, dass für jedes der unendlich vielen Präfixe
$bab,babab,bababab,\dots$ von $(ba)^\omega$ das entsprechende Präfix von $r$
zwar den akzeptierenden Zustand $\{q_0,q_1\}$ erreicht,
aber der zugehörige in $q_0$ endende Teilrun von $\Amc$
nicht mehr zu einem erfolgreichen Run auf $\alpha$
fortgesetzt werden kann.

\goodbreak
% ===================================================================
\section*{T3.13~ Variation der Akzeptanzbedingung im vorigen Beispiel}

Wir betrachten dieselben Automaten $\Amc,\Amc^d$ wie im vorigen Beispiel.
Man könnte versuchen, durch Variation der Akzeptanzbedingung von $\Amc^d$
einen zu $\Amc$ äquivalenten deterministischen Muller-, Rabin- oder Streett-Automaten
zu erhalten, ohne die eigentliche Potenzmengenkonstruktion aufzugeben.
Dieser Versuch muss aber scheitern, wovon man sich leicht überzeugt,
wenn man alle möglichen Akzeptanzbedingungen systematisch durchgeht.
Diese sind entweder trivial (d.\,h.\ führen offensichtlich zu $L_\omega(\Amc^d) = \emptyset$
oder $L_\omega(\Amc^d) = \Sigma^\omega$)
oder laufen auf einen der folgenden Fälle hinaus:
Erfolgreiche Runs \dots
%
\begin{enumerate}
  \item
    \dots\ müssen $\{q_1\}$ und $\{q_0,q_1\}$ unendlich oft besuchen;
  \item
    \dots\ dürfen nur $\{q_1\}$ unendlich oft besuchen;
  \item
    \dots\ dürfen nur $\{q_0,q_1\}$ unendlich oft besuchen.
\end{enumerate}
%
Im 1.\ und 2.\ Fall gibt es jedoch Bad Runs auf $(ba)^\omega$ bzw.\ $ba^\omega$;
im 3.\ Fall gibt es Wörter, die von $\Amc$ akzeptiert werden,
aber nicht von $\Amc^d$, z.\,B.\ $(bba)^\omega$.

% ===================================================================
\section*{T3.14~ Beispiel für Safras Trick 1}

Betrachte folgenden NBA $\Amc$ über dem Alphabet $\Sigma = \{a,b\}$.
%
\begin{center}
  \begin{tikzpicture}[%
    node distance=25mm,>=Latex,
    initial text="", initial where=below left,
    every state/.style={draw=black,thin,fill=black!5,inner sep=1mm,minimum size=6mm},
    accepting/.style={double distance=1.5pt, double=white},
    every edge/.style={draw=black,thin}
  ]
    \node[state,initial]   (q0) {$q_0$};
    \node[state,accepting] (q1) [right of=q0] {$q_1$};
    \node[above left=-2mm and 8mm of q0] {$\Amc$};
    
    \path[->] (q0) edge              node [above] {$b$}    (q1)
              (q0) edge [loop above] node [right] {~$a,b$} ()
              (q1) edge [loop above] node [right] {~$a$}   ();
  \end{tikzpicture}
\end{center}
%
Dieser NBA akzeptiert genau die $\omega$-Wörter mit endlich vielen $b$'s.
Die Potenzmengen\-konstruktion liefert den DBA
%
\begin{center}
  \begin{tikzpicture}[%
    node distance=40mm,>=Latex,
    initial text="", initial where=below left,
    every state/.style={ellipse,draw=black,thin,fill=black!5,inner sep=1mm,minimum size=6mm},
    accepting/.style={double distance=1.5pt, double=white},
    every edge/.style={draw=black,thin}
  ]
    \node[state,initial]   (q0)                      {$\{q_0\}$};
    \node[state,accepting] (q0q1) [right of=q0]      {$\{q_0,q_1\}$};
%    \node[above left=-2mm and 8mm of q0q21] {$\Amc^d$};
    
    \path[->] (q0)   edge [bend left=10]  node [above] {$a$}  (q0q1)
              (q0q1) edge [bend left=10]  node [below] {$b$}  (q0)
              (q0)   edge [loop above]    node [right] {~$b$} ()
              (q0q1) edge [loop above]    node [right] {~$a$} ();
  \end{tikzpicture}
\end{center}
%
mit einem Bad Run auf $(ab)^\omega$.
Mittels Safras Trick 1 erhält man hingegen folgenden deterministischen Automaten $\Amc^d$.
%
\begin{center}
  \begin{tikzpicture}[%
    node distance=40mm,>=Latex,
    initial text="", initial where=below left,
    every state/.style={ellipse,draw=black,thin,fill=black!5,inner sep=.7mm},
    accepting/.style={double distance=1.5pt, double=white},
    every edge/.style={draw=black,thin}
  ]
    \begin{scope}[
      every node/.style={rectangle,rounded corners,draw=black,fill=white,inner sep=1mm}
    ]    
      \node (S0eps)                      {$\{q_0\}$};
      \node (S1eps) [right of=S0eps]     {$\{q_0,q_1\}$};
      \node (S2eps) [right of=S1eps]     {$\{q_0,q_1\}$};
      \node (S20)   [below=4mm of S2eps] {$\{q_1\}$};
    \end{scope}

    \begin{scope}[
      every node/.style={rectangle,rounded corners,draw=none,fill=none,inner sep=1mm}
    ]    
      \node (S00)   [below=4mm of S0eps] {$\phantom{\{q_1\}}$};
      \node (S10)   [below=4mm of S1eps] {$\phantom{\{q_1\}}$};
    \end{scope}
  
    \begin{pgfonlayer}{background}
      \node[state,initial] (S0) [fit=(S0eps) (S00)] {};
      \node[state]         (S1) [fit=(S1eps) (S10)] {};
      \node[state]         (S2) [fit=(S2eps) (S20)] {};
  %    \node[above left=-2mm and 8mm of q0q21] {$\Amc^d$};
    \end{pgfonlayer}  
    
%    \node [above right=-1mm and -1mm of S0] {$S_0$};
%    \node [above right=-1mm and -1mm of S1] {$S_1$};
%    \node [above right=-1mm and -1mm of S2] {$S_2$};
%
    \node [below right=-1mm and -1mm of S0] {$S_0$};
    \node [below right=-1mm and -1mm of S1] {$S_1$};
    \node [below right=-1mm and -1mm of S2] {$S_2$};

%    \node [below left=-1mm   and 1mm of S0.east] {$S_0$};
%    \node [below left=-1.5mm and .5mm of S1.east] {$S_1$};
%    \node [below left=-2mm   and 0mm of S2.east] {$S_2$};
    
    \node [above left=-2mm and 8mm of q0q21] {$\Amc^d$};
    
    \path[->] (S0) edge [bend left=10]      node [above] {$a$}  (S1)
              (S1) edge [bend left=10]      node [below] {$b$}  (S0)
              (S1) edge                     node [above] {$a$}  (S2)
              (S2) edge [bend left=35]      node [very near start,above] {$b$}  (S0)
              (S0) edge [out=98,in=82,loop] node [right] {~$b$} ()
              (S2) edge [out=98,in=82,loop] node [right] {~$a$} ();

    \path[-]  (S2eps) edge (S20);
  \end{tikzpicture}
\end{center}
%
Wenn man nun den Bad Run auf dem Wort $(ab)^\omega$ verhindern möchte,
dann muss man die Akzeptanzbedingung so wählen, dass
$S_2$ unendlich oft besucht werden muss, aber $S_0$ und $S_1$ nur
endlich oft. Dies erreicht man z.\,B.\ durch die Rabin-Akzeptanzkomponente
$\Pmc = \{(\{S_0\},\{S_2\})\}$.

% ===================================================================
\section*{T3.15~ Beispiel für die gesamte Safra-Konstruktion}

Wir betrachten den NBA $\Amc$ aus dem vorigen Beispiel:
%
\begin{center}
  \begin{tikzpicture}[%
    node distance=25mm,>=Latex,
    initial text="", initial where=below left,
    every state/.style={draw=black,thin,fill=black!5,inner sep=1mm,minimum size=6mm},
    accepting/.style={double distance=1.5pt, double=white},
    every edge/.style={draw=black,thin}
  ]
    \node[state,initial]   (q0) {$q_0$};
    \node[state,accepting] (q1) [right of=q0] {$q_1$};
    \node[above left=-2mm and 8mm of q0] {$\Amc$};
    
    \path[->] (q0) edge              node [above] {$b$}    (q1)
              (q0) edge [loop above] node [right] {~$a,b$} ()
              (q1) edge [loop above] node [right] {~$a$}   ();
  \end{tikzpicture}
\end{center}
%
Im Folgenden wird die Konstruktion des DRA $\Amc^d$ gemäß der Safra-Konstruktion
schrittweise beschrieben. Dabei benennen wir die konstruierten Zustände
(Safrabäume) der Reihe nach mit $S_0,S_1,S_2,\dots$ und schreiben diese Namen
jeweils rechts neben den entsprechenden Zustand.
Außerdem verwenden wir innerhalb von Safrabäumen als Knotennamen
die Zahlen $1,2,3,\dots$ und schreiben sie
rechts neben den jeweiligen Knoten. Markierte Knoten (Schritt 6) werden wie gehabt
mit \circled{!} gekennzeichnet.

\paragraph*{Startzustand} ist der folgende Safrabaum:
%
\begin{center}
  \begin{tikzpicture}[%
    node distance=40mm,>=Latex,
    initial text="", initial where=below left,
    every state/.style={ellipse,draw=black,thin,fill=black!5,inner sep=.7mm},
    accepting/.style={double distance=1.5pt, double=white},
    every edge/.style={draw=black,thin}
  ]
    \begin{scope}[
      every node/.style={rectangle,rounded corners,draw=black,fill=white,inner sep=1mm}
    ]    
      \node (S01)                      {$\{q_0\}$};
    \end{scope}

    \node (S01l)  [right=0mm of S01] {1};
  
    \begin{pgfonlayer}{background}
      \node[state,initial] (S0) [fit=(S01) (S01l)] {};
    \end{pgfonlayer}  
    
    \node [below right=-1mm and -1mm of S0] {$S_0$};
  \end{tikzpicture}
\end{center}

\paragraph*{{\boldmath Folgezustand von $S_0$ mit $b$}}
%
\begin{itemize}
  \item
    Schritt 1 der Safra-Konstruktion ist nicht anwendbar,
    da Knoten 1 nicht markiert ist.
  \item
    Schritt 2 ist nicht anwendbar, da Knoten 1 keine akzeptierenden Zustände enthält
    ($F = \{q_1\}$, siehe Bild).
  \item
    In Schritt 3 ändert sich der Makrozustand von Knoten 1 nicht, 
    da der einzige $b$-Nachfolgezustand von $q_0$ wieder $q_0$ ist.
  \item
    Schritte 4--6 sind nicht anwendbar, da Knoten 1 noch keine Kinder hat.
\end{itemize}
Also ist der $b$-Folgezustand von $S_0$ wieder $S_0$:
%
\begin{center}
  \begin{tikzpicture}[%
    node distance=40mm,>=Latex,
    initial text="", initial where=below left,
    every state/.style={ellipse,draw=black,thin,fill=black!5,inner sep=.7mm},
    accepting/.style={double distance=1.5pt, double=white},
    every edge/.style={draw=black,thin}
  ]
    \begin{scope}[
      every node/.style={rectangle,rounded corners,draw=black,fill=white,inner sep=1mm}
    ]    
      \node (S01)                      {$\{q_0\}$};
    \end{scope}

    \node (S01l)  [right=0mm of S01] {1};
  
    \begin{pgfonlayer}{background}
      \node[state,initial] (S0) [fit=(S01) (S01l)] {};
    \end{pgfonlayer}  
    
    \node [below right=-1mm and -1mm of S0] {$S_0$};

    \path[->] (S0) edge [loop above] node [right] {~$b$} ();
  \end{tikzpicture}
\end{center}

\paragraph*{{\boldmath Folgezustand von $S_0$ mit $a$}}
%
\begin{itemize}
  \item
    Schritte 1--2 sind nicht anwendbar, siehe oben.
  \item
    In Schritt 3 ändert sich der Makrozustand von Knoten 1 
    zu $\{q_0,q_1\}$, da sowohl $q_0$ als auch $q_1$
    von $q_0$ aus mit $a$ erreicht werden können.
  \item
    Schritte 4--6 sind nicht anwendbar, da Knoten 1 noch keine Kinder hat.
\end{itemize}
%
Also ist der $a$-Folgezustand von $S_0$ ein neuer Safrabaum $S_1$,
in dem Knoten 1 den Makrozustand $\{q_0,q_1\}$ hat:
%
\begin{center}
  \begin{tikzpicture}[%
    node distance=40mm,>=Latex,
    initial text="", initial where=below left,
    every state/.style={ellipse,draw=black,thin,fill=black!5,inner sep=.7mm},
    accepting/.style={double distance=1.5pt, double=white},
    every edge/.style={draw=black,thin}
  ]
    \begin{scope}[
      every node/.style={rectangle,rounded corners,draw=black,fill=white,inner sep=1mm}
    ]    
      \node (S01)                      {$\{q_0\}$};
      \node (S11) [right=of S01]       {$\{q_0,q_1\}$};
    \end{scope}

    \node (S01l)  [right=0mm of S01] {1};
    \node (S11l)  [right=0mm of S11] {1};
  
    \begin{pgfonlayer}{background}
      \node[state,initial] (S0) [fit=(S01) (S01l)] {};
      \node[state]         (S1) [fit=(S11) (S11l)] {};
    \end{pgfonlayer}  
    
    \node [below right=-1mm and -1mm of S0] {$S_0$};
    \node [below right=-1mm and -1mm of S1] {$S_1$};

    \path[->] (S0) edge              node [above] {$a$}  (S1)
              (S0) edge [loop above] node [right] {~$b$} ();
  \end{tikzpicture}
\end{center}

\paragraph*{{\boldmath Folgezustand von $S_1$ mit $a$}}
%
\begin{itemize}
  \item
    Schritt 1 ist nach wie vor nicht anwendbar (keine Markierung).
  \item
    In Schritt 2 wird ein neues Kind von Knoten 1 erzeugt,
    dessen Makrozustand aus dem akzeptierenden Zustand $q_1$
    aus Knoten 1 besteht und das den Namen 2 bekommt:
    %    
    \begin{center}
      \begin{tikzpicture}[%
        node distance=40mm,>=Latex,
        initial text="", initial where=below left,
        every state/.style={ellipse,draw=black,thin,fill=black!5,inner sep=.7mm},
        accepting/.style={double distance=1.5pt, double=white},
        every edge/.style={draw=black,thin}
      ]
        \begin{scope}[
          every node/.style={rectangle,rounded corners,draw=black,fill=white,inner sep=1mm}
        ]    
          \node (S11)                    {$\{q_0,q_1\}$};
          \node (S12) [below=4mm of S11] {$\{q_1\}$};
        \end{scope}
    
        \node (S11l)  [right=0mm of S11] {1};
        \node (S12l)  [right=0mm of S12] {2};
      
        \path[-] (S11) edge (S12);
      \end{tikzpicture}
    \end{center}
  \item
    In Schritt 3 wird auf beide Knoten die Potenzmengenkonstruktion angewendet;
    bei Übergang mit $a$ ändert sich der Inhalt beider Makrozustände nicht.
  \item
    Schritt 4 ist nicht anwendbar, da kein Knoten mehr als ein Kind hat.
  \item
    Schritt 5 ist nicht anwendbar, da der Makrozustand von Knoten 2 nicht leer ist.
  \item
    Schritt 6 ist nicht anwendbar, da $q_0$ im Makrozustand von Knoten 1,
    aber nicht von Knoten 2 vorkommt. 
\end{itemize}
%
Also ist der $a$-Folgezustand von $S_1$ ein neuer Safrabaum $S_2$
mit zwei Knoten wie folgt:
%
\begin{center}
  \begin{tikzpicture}[%
    node distance=40mm,>=Latex,
    initial text="", initial where=below left,
    every state/.style={ellipse,draw=black,thin,fill=black!5,inner sep=.7mm},
    accepting/.style={double distance=1.5pt, double=white},
    every edge/.style={draw=black,thin}
  ]
    \begin{scope}[
      every node/.style={rectangle,rounded corners,draw=black,fill=white,inner sep=1mm}
    ]    
      \node (S01)                      {$\{q_0\}$};
      \node (S11) [right=of S01]       {$\{q_0,q_1\}$};
      \node (S21) [below=26mm of S11]  {$\{q_0,q_1\}$};
      \node (S22) [below=4mm of S21]   {$\{q_1\}$};
    \end{scope}

    \node (S01l)  [right=0mm of S01] {1};
    \node (S11l)  [right=0mm of S11] {1};
    \node (S21l)  [right=0mm of S21] {1};
    \node (S22l)  [right=0mm of S22] {2};
  
    \begin{scope}[
      every node/.style={rectangle,rounded corners,draw=none,fill=none,inner sep=1mm}
    ]    
      \node (S02)   [below=4mm of S01] {$\phantom{\{q_1\}}$};
      \node (S12)   [below=4mm of S11] {$\phantom{\{q_1\}}$};
    \end{scope}
  
    \begin{pgfonlayer}{background}
      \node[state,initial] (S0) [fit=(S01) (S01l) (S02)] {};
      \node[state]         (S1) [fit=(S11) (S11l) (S12)] {};
      \node[state]         (S2) [fit=(S21) (S21l) (S22) (S22l)] {};
    \end{pgfonlayer}  
    
    \node [below right=-1mm and -1mm of S0] {$S_0$};
    \node [below right=-1mm and -1mm of S1] {$S_1$};
    \node [below right=-1mm and -1mm of S2] {$S_2$};

    \path[->] (S0) edge                     node [above] {$a$}  (S1)
              (S1) edge                     node [right] {$a$}  (S2)
              (S0) edge [out=98,in=82,loop] node [right] {~$b$} ();
              
    \path[-] (S21) edge (S22);
  \end{tikzpicture}
\end{center}

\paragraph*{{\boldmath Folgezustand von $S_1$ mit $b$}}
%
\begin{itemize}
  \item
    Schritte 1--2 wie oben.
  \item
    In Schritt 3 wird wieder auf beide Knoten die Potenzmengenkonstruktion
    angewendet; bei Übergang mit $b$ ändern sich die Makrozustände wie folgt:
    %    
    \begin{center}
      \begin{tikzpicture}[%
        node distance=40mm,>=Latex,
        initial text="", initial where=below left,
        every state/.style={ellipse,draw=black,thin,fill=black!5,inner sep=.7mm},
        accepting/.style={double distance=1.5pt, double=white},
        every edge/.style={draw=black,thin}
      ]
        \begin{scope}[
          every node/.style={rectangle,rounded corners,draw=black,fill=white,inner sep=1mm}
        ]    
          \node (S11)                    {$\{q_0\}$};
          \node (S12) [below=4mm of S11] {$\emptyset$};
        \end{scope}
    
        \node (S11l)  [right=0mm of S11] {1};
        \node (S12l)  [right=0mm of S12] {2};
      
        \path[-] (S11) edge (S12);
      \end{tikzpicture}
    \end{center}
  \item
    Schritt 4 ist nicht anwendbar, da kein Knoten mehr als ein Kind hat.
  \item
    In Schritt 5 wird Knoten 2 gelöscht.
  \item
    Schritt 6 ist nicht anwendbar, da Knoten 1 nun kein Kind mehr hat.
\end{itemize}
%
Damit ist der $b$-Folgezustand von $S_1$ der Safrabaum $S_0$:
%
\begin{center}
  \begin{tikzpicture}[%
    node distance=40mm,>=Latex,
    initial text="", initial where=below left,
    every state/.style={ellipse,draw=black,thin,fill=black!5,inner sep=.7mm},
    accepting/.style={double distance=1.5pt, double=white},
    every edge/.style={draw=black,thin}
  ]
    \begin{scope}[
      every node/.style={rectangle,rounded corners,draw=black,fill=white,inner sep=1mm}
    ]    
      \node (S01)                      {$\{q_0\}$};
      \node (S11) [right=of S01]       {$\{q_0,q_1\}$};
      \node (S21) [below=26mm of S11]  {$\{q_0,q_1\}$};
      \node (S22) [below=4mm of S21]   {$\{q_1\}$};
    \end{scope}

    \node (S01l)  [right=0mm of S01] {1};
    \node (S11l)  [right=0mm of S11] {1};
    \node (S21l)  [right=0mm of S21] {1};
    \node (S22l)  [right=0mm of S22] {2};
  
    \begin{scope}[
      every node/.style={rectangle,rounded corners,draw=none,fill=none,inner sep=1mm}
    ]    
      \node (S02)   [below=4mm of S01] {$\phantom{\{q_1\}}$};
      \node (S12)   [below=4mm of S11] {$\phantom{\{q_1\}}$};
    \end{scope}
  
    \begin{pgfonlayer}{background}
      \node[state,initial] (S0) [fit=(S01) (S01l) (S02)] {};
      \node[state]         (S1) [fit=(S11) (S11l) (S12)] {};
      \node[state]         (S2) [fit=(S21) (S21l) (S22) (S22l)] {};
    \end{pgfonlayer}  
    
    \node [below right=-1mm and -1mm of S0] {$S_0$};
    \node [below right=-1mm and -1mm of S1] {$S_1$};
    \node [below right=-1mm and -1mm of S2] {$S_2$};

    \path[->] (S0) edge [bend left=10]      node [above] {$a$}  (S1)
              (S1) edge [bend left=10]      node [below] {$b$}  (S0)
              (S1) edge                     node [right] {$a$}  (S2)
              (S0) edge [out=98,in=82,loop] node [right] {~$b$} ();
              
    \path[-] (S21) edge (S22);
  \end{tikzpicture}
\end{center}

\paragraph*{{\boldmath Folgezustand von $S_2$ mit $a$}}
%
\begin{itemize}
  \item
    Schritt 1 ist nach wie vor nicht anwendbar.
  \item
    In Schritt 2 wird je ein neues Kind von Knoten 1 und 2 erzeugt,
    da die Makrozustände beider Knoten akzeptierende Zustände enthalten:
    %    
    \begin{center}
      \begin{tikzpicture}[%
        node distance=40mm,>=Latex,
        initial text="", initial where=below left,
        every state/.style={ellipse,draw=black,thin,fill=black!5,inner sep=.7mm},
        accepting/.style={double distance=1.5pt, double=white},
        every edge/.style={draw=black,thin}
      ]
        \begin{scope}[
          every node/.style={rectangle,rounded corners,draw=black,fill=white,inner sep=1mm}
        ]    
          \node (S11)                                   {$\{q_0,q_1\}$};
          \node (S12) [below left =4mm and -2mm of S11] {$\{q_1\}$};
          \node (S13) [below right=4mm and -2mm of S11] {$\{q_1\}$};
          \node (S14) [below=4mm of S12]                {$\{q_1\}$};
        \end{scope}
    
        \node (S11l)  [right=0mm of S11] {1};
        \node (S12l)  [right=0mm of S12] {2};
        \node (S13l)  [right=0mm of S13] {3};
        \node (S14l)  [right=0mm of S14] {4};
      
        \path[-] (S11) edge (S12)
                 (S11) edge (S13)
                 (S12) edge (S14);
      \end{tikzpicture}
    \end{center}
  \item
    In Schritt 3 wird auf alle Knoten die Potenzmengenkonstruktion angewendet;
    bei Übergang mit $a$ ändert sich der Inhalt der Makrozustände nicht.
  \item
    In Schritt 4 wird $q_1$ aus Knoten 3 entfernt, da dieser Zustand im älteren
    Geschwister 2 enthalten ist:
    %    
    \begin{center}
      \begin{tikzpicture}[%
        node distance=40mm,>=Latex,
        initial text="", initial where=below left,
        every state/.style={ellipse,draw=black,thin,fill=black!5,inner sep=.7mm},
        accepting/.style={double distance=1.5pt, double=white},
        every edge/.style={draw=black,thin}
      ]
        \begin{scope}[
          every node/.style={rectangle,rounded corners,draw=black,fill=white,inner sep=1mm}
        ]    
          \node (S11)                                   {$\{q_0,q_1\}$};
          \node (S12) [below left =4mm and -2mm of S11] {$\{q_1\}$};
          \node (S13) [below right=4mm and -2mm of S11] {$\emptyset$};
          \node (S14) [below=4mm of S12]                {$\{q_1\}$};
        \end{scope}
    
        \node (S11l)  [right=0mm of S11] {1};
        \node (S12l)  [right=0mm of S12] {2};
        \node (S13l)  [right=0mm of S13] {3};
        \node (S14l)  [right=0mm of S14] {4};
      
        \path[-] (S11) edge (S12)
                 (S11) edge (S13)
                 (S12) edge (S14);
      \end{tikzpicture}
    \end{center}
  \item
    In Schritt 5 wird Knoten 3 gelöscht:
    %    
    \begin{center}
      \begin{tikzpicture}[%
        node distance=40mm,>=Latex,
        initial text="", initial where=below left,
        every state/.style={ellipse,draw=black,thin,fill=black!5,inner sep=.7mm},
        accepting/.style={double distance=1.5pt, double=white},
        every edge/.style={draw=black,thin}
      ]
        \begin{scope}[
          every node/.style={rectangle,rounded corners,draw=black,fill=white,inner sep=1mm}
        ]    
          \node (S11)                                   {$\{q_0,q_1\}$};
          \node (S12) [below left =4mm and -2mm of S11] {$\{q_1\}$};
          \node (S14) [below=4mm of S12]                {$\{q_1\}$};
        \end{scope}
    
        \node (S11l)  [right=0mm of S11] {1};
        \node (S12l)  [right=0mm of S12] {2};
        \node (S14l)  [right=0mm of S14] {4};
      
        \path[-] (S11) edge (S12)
                 (S12) edge (S14);
      \end{tikzpicture}
    \end{center}
  \item
    In Schritt 6 wird nun Knoten 4 gelöscht und Knoten 2 markiert:
    %    
    \begin{center}
      \begin{tikzpicture}[%
        node distance=40mm,>=Latex,
        initial text="", initial where=below left,
        every state/.style={ellipse,draw=black,thin,fill=black!5,inner sep=.7mm},
        accepting/.style={double distance=1.5pt, double=white},
        every edge/.style={draw=black,thin}
      ]
        \begin{scope}[
          every node/.style={rectangle,rounded corners,draw=black,fill=white,inner sep=1mm}
        ]    
          \node (S11)                                   {$\{q_0,q_1\}$};
          \node (S12) [below left=4mm and -2mm of S11] {$\{q_1\}$};
        \end{scope}
    
        \node (S11l)  [right=0mm of S11] {1};
        \node (S12l)  [right=0mm of S12] {2};
        \node (S12m)  [left =0mm of S12] {\circled{!}};
      
        \path[-] (S11) edge (S12);
      \end{tikzpicture}
    \end{center}
\end{itemize}
%
Der zuletzt abgebildete Safrabaum $S_3$ ist der $a$-Folgezustand von $S_2$:
%
\begin{center}
  \begin{tikzpicture}[%
    node distance=40mm,>=Latex,
    initial text="", initial where=below left,
    every state/.style={ellipse,draw=black,thin,fill=black!5,inner sep=.7mm},
    accepting/.style={double distance=1.5pt, double=white},
    every edge/.style={draw=black,thin}
  ]
    \begin{scope}[
      every node/.style={rectangle,rounded corners,draw=black,fill=white,inner sep=1mm}
    ]    
      \node (S01)                      {$\{q_0\}$};
      \node (S11) [right=of S01]       {$\{q_0,q_1\}$};
      \node (S21) [below=26mm of S11]  {$\{q_0,q_1\}$};
      \node (S22) [below=4mm of S21]   {$\{q_1\}$};
      \node (S31) [left=of S21]        {$\{q_0,q_1\}$};
      \node (S32) [below=4mm of S31]   {$\{q_1\}$};
    \end{scope}

    \node (S01l)  [right=0mm of S01] {1};
    \node (S11l)  [right=0mm of S11] {1};
    \node (S21l)  [right=0mm of S21] {1};
    \node (S22l)  [right=0mm of S22] {2};
    \node (S31l)  [right=0mm of S31] {1};
    \node (S32l)  [right=0mm of S32] {2};
    \node (S32m)  [left =0mm of S32] {\circled{!}};
  
    \begin{scope}[
      every node/.style={rectangle,rounded corners,draw=none,fill=none,inner sep=1mm}
    ]    
      \node (S02)   [below=4mm of S01] {$\phantom{\{q_1\}}$};
      \node (S12)   [below=4mm of S11] {$\phantom{\{q_1\}}$};
    \end{scope}
  
    \begin{pgfonlayer}{background}
      \node[state,initial] (S0) [fit=(S01) (S01l) (S02)] {};
      \node[state]         (S1) [fit=(S11) (S11l) (S12)] {};
      \node[state]         (S2) [fit=(S21) (S21l) (S22) (S22l)] {};
      \node[state]         (S3) [fit=(S31) (S31l) (S32) (S32l) (S32m)] {};
    \end{pgfonlayer}  
    
    \node [below right=-1mm and -1mm of S0] {$S_0$};
    \node [below right=-1mm and -1mm of S1] {$S_1$};
    \node [below right=-1mm and -1mm of S2] {$S_2$};
    \node [below right=-1mm and -1mm of S3] {$S_3$};

    \path[->] (S0) edge [bend left=10]      node [above] {$a$}  (S1)
              (S1) edge [bend left=10]      node [below] {$b$}  (S0)
              (S1) edge                     node [right] {$a$}  (S2)
              (S0) edge [out=98,in=82,loop] node [right] {~$b$} ()
              (S2) edge                     node [above] {$a$}  (S3);
              
    \path[-] (S21) edge (S22)
             (S31) edge (S32);
  \end{tikzpicture}
\end{center}

\paragraph*{{\boldmath Folgezustand von $S_2$ mit $b$}}
%
\begin{itemize}
  \item
    Schritte 1--2 wie oben.
  \item
    In Schritt 3 wird wieder auf alle Knoten die Potenzmengenkonstruktion
    angewendet; bei Übergang mit $b$ ändern sich die Makrozustände wie folgt:
    %    
    \begin{center}
      \begin{tikzpicture}[%
        node distance=40mm,>=Latex,
        initial text="", initial where=below left,
        every state/.style={ellipse,draw=black,thin,fill=black!5,inner sep=.7mm},
        accepting/.style={double distance=1.5pt, double=white},
        every edge/.style={draw=black,thin}
      ]
        \begin{scope}[
          every node/.style={rectangle,rounded corners,draw=black,fill=white,inner sep=1mm}
        ]    
          \node (S11)                                   {$\{q_0,q_1\}$};
          \node (S12) [below left =4mm and -2mm of S11] {$\emptyset$};
          \node (S13) [below right=4mm and -2mm of S11] {$\emptyset$};
          \node (S14) [below=4mm of S12]                {$\emptyset$};
        \end{scope}
    
        \node (S11l)  [right=0mm of S11] {1};
        \node (S12l)  [right=0mm of S12] {2};
        \node (S13l)  [right=0mm of S13] {3};
        \node (S14l)  [right=0mm of S14] {4};
      
        \path[-] (S11) edge (S12)
                 (S11) edge (S13)
                 (S12) edge (S14);
      \end{tikzpicture}
    \end{center}
  \item
    Schritt 4 ist nicht anwendbar, da alle Makrozsutände außer dem von Knoten 1 leer sind.
  \item
    In Schritt 5 werden Knoten $2,3,4$ gelöscht.
  \item
    Schritt 6 ist nicht anwendbar, da Knoten 1 nun kein Kind mehr hat.
\end{itemize}
%
\enlargethispage*{20mm}
Damit ist der $b$-Folgezustand von $S_1$ wieder der Safrabaum $S_0$:
%
\begin{center}
  \begin{tikzpicture}[%
    node distance=40mm,>=Latex,
    initial text="", initial where=below left,
    every state/.style={ellipse,draw=black,thin,fill=black!5,inner sep=.7mm},
    accepting/.style={double distance=1.5pt, double=white},
    every edge/.style={draw=black,thin}
  ]
    \begin{scope}[
      every node/.style={rectangle,rounded corners,draw=black,fill=white,inner sep=1mm}
    ]    
      \node (S01)                      {$\{q_0\}$};
      \node (S11) [right=of S01]       {$\{q_0,q_1\}$};
      \node (S21) [below=26mm of S11]  {$\{q_0,q_1\}$};
      \node (S22) [below=4mm of S21]   {$\{q_1\}$};
      \node (S31) [left=of S21]        {$\{q_0,q_1\}$};
      \node (S32) [below=4mm of S31]   {$\{q_1\}$};
    \end{scope}

    \node (S01l)  [right=0mm of S01] {1};
    \node (S11l)  [right=0mm of S11] {1};
    \node (S21l)  [right=0mm of S21] {1};
    \node (S22l)  [right=0mm of S22] {2};
    \node (S31l)  [right=0mm of S31] {1};
    \node (S32l)  [right=0mm of S32] {2};
    \node (S32m)  [left =0mm of S32] {\circled{!}};
  
    \begin{scope}[
      every node/.style={rectangle,rounded corners,draw=none,fill=none,inner sep=1mm}
    ]    
      \node (S02)   [below=4mm of S01] {$\phantom{\{q_1\}}$};
      \node (S12)   [below=4mm of S11] {$\phantom{\{q_1\}}$};
    \end{scope}
  
    \begin{pgfonlayer}{background}
      \node[state,initial] (S0) [fit=(S01) (S01l) (S02)] {};
      \node[state]         (S1) [fit=(S11) (S11l) (S12)] {};
      \node[state]         (S2) [fit=(S21) (S21l) (S22) (S22l)] {};
      \node[state]         (S3) [fit=(S31) (S31l) (S32) (S32l) (S32m)] {};
    \end{pgfonlayer}  
    
    \node [below right=-1mm and -1mm of S0] {$S_0$};
    \node [below right=-1mm and -1mm of S1] {$S_1$};
    \node [below right=-1mm and -1mm of S2] {$S_2$};
    \node [below right=-1mm and -1mm of S3] {$S_3$};

    \path[->] (S0) edge [bend left=10]      node [above] {$a$}  (S1)
              (S1) edge [bend left=10]      node [below] {$b$}  (S0)
              (S1) edge                     node [right] {$a$}  (S2)
              (S0) edge [out=98,in=82,loop] node [right] {~$b$} ()
              (S2) edge                     node [above] {$a$}  (S3)
              (S2) edge                     node [above,near start] {$b$}  (S0);
              
    \path[-] (S21) edge (S22)
             (S31) edge (S32);
  \end{tikzpicture}
\end{center}

\paragraph*{{\boldmath Folgezustände von $S_3$}}
%
~\par
Da sich der Safrabaum $S_3$ von $S_2$ nur durch die Markierung des Knotens 2 unterscheidet,
laufen die Schritte 2--6 genauso ab, nachdem in Schritt 1 die Markierung entfernt wurde.
Damit hat $S_3$ dieselben Folgezustände wie $S_2$ (nämlich $S_3$ für $a$ und $S_1$ für $b$):
%
\begin{center}
  \begin{tikzpicture}[%
    node distance=40mm,>=Latex,
    initial text="", initial where=below left,
    every state/.style={ellipse,draw=black,thin,fill=black!5,inner sep=.7mm},
    accepting/.style={double distance=1.5pt, double=white},
    every edge/.style={draw=black,thin}
  ]
    \begin{scope}[
      every node/.style={rectangle,rounded corners,draw=black,fill=white,inner sep=1mm}
    ]    
      \node (S01)                      {$\{q_0\}$};
      \node (S11) [right=of S01]       {$\{q_0,q_1\}$};
      \node (S21) [below=26mm of S11]  {$\{q_0,q_1\}$};
      \node (S22) [below=4mm of S21]   {$\{q_1\}$};
      \node (S31) [left=of S21]        {$\{q_0,q_1\}$};
      \node (S32) [below=4mm of S31]   {$\{q_1\}$};
    \end{scope}

    \node (S01l)  [right=0mm of S01] {1};
    \node (S11l)  [right=0mm of S11] {1};
    \node (S21l)  [right=0mm of S21] {1};
    \node (S22l)  [right=0mm of S22] {2};
    \node (S31l)  [right=0mm of S31] {1};
    \node (S32l)  [right=0mm of S32] {2};
    \node (S32m)  [left =0mm of S32] {\circled{!}};
  
    \begin{scope}[
      every node/.style={rectangle,rounded corners,draw=none,fill=none,inner sep=1mm}
    ]    
      \node (S02)   [below=4mm of S01] {$\phantom{\{q_1\}}$};
      \node (S12)   [below=4mm of S11] {$\phantom{\{q_1\}}$};
    \end{scope}
  
    \begin{pgfonlayer}{background}
      \node[state,initial] (S0) [fit=(S01) (S01l) (S02)] {};
      \node[state]         (S1) [fit=(S11) (S11l) (S12)] {};
      \node[state]         (S2) [fit=(S21) (S21l) (S22) (S22l)] {};
      \node[state]         (S3) [fit=(S31) (S31l) (S32) (S32l) (S32m)] {};
    \end{pgfonlayer}  
    
    \node [below right=-1mm and -1mm of S0] {$S_0$};
    \node [below right=-1mm and -1mm of S1] {$S_1$};
    \node [below right=-1mm and -1mm of S2] {$S_2$};
    \node [below right=-1mm and -1mm of S3] {$S_3$};

    \path[->] (S0) edge [bend left=10]        node [above] {$a$}  (S1)
              (S1) edge [bend left=10]        node [below] {$b$}  (S0)
              (S1) edge                       node [right] {$a$}  (S2)
              (S0) edge [out=98,in=82,loop]   node [right] {~$b$} ()
              (S2) edge                       node [above] {$a$}  (S3)
              (S2) edge                       node [above,near start] {$b$} (S0)
              (S3) edge                       node [left]  {$a$}  (S0)
              (S3) edge [out=186,in=174,loop] node [left] {~$b$} ();
              
    \path[-] (S21) edge (S22)
             (S31) edge (S32);
  \end{tikzpicture}
\end{center}
%
Damit sind alle erreichbaren Zustände (Safrabäume) erzeugt.

\paragraph*{{\boldmath Akzeptanzkomponente}}
%
~\par
Laut Folie~72 ist $\Pmc = \{(E_1,F_1),\,(E_2,F_2)\}
= \{(\emptyset,\emptyset),\,(\{S_1,S_1\},\{S_3\})\}$.

% ===================================================================
\section*{T3.16~ Korrektheitsbeweis der Safra-Konstruktion, Details}

\textsfbf{Hilfsaussage [HA].}~
Für alle $T_i$ und alle Zustände $p$ im Makrozustand \textsfbf{(MZ)} von $v$ in $T_{i+1}$
gibt es einen Zustand $q$ im Makrozustand von $v$ in $T_{i}$
und einen endlichen Run $q\dots p$ von $\Amc$ auf dem zugehörigen Teilwort von $\alpha$,
der einen akzeptierenden Zustand enthält.

\par\medskip\noindent
\textsfbf{Skizze:}~
%
\begin{center}
  \begin{tikzpicture}[%
    node distance=40mm,>=Latex,
    initial text="", initial where=below left,
    every state/.style={ellipse,draw=black,thin,fill=black!5,inner sep=.7mm},
    accepting/.style={double distance=1.5pt, double=white},
    every edge/.style={draw=black,thin}
  ]
    \begin{scope}[
      every node/.style={rectangle,rounded corners,draw=black,fill=white,inner sep=1mm}
    ]    
      \node (vi)                {\strut$\cdots q\cdots$};
      \node (vi1) [right=of vi] {\strut$\cdots p\cdots$};
    \end{scope}

    \node (vil)  [above right=-3.5mm and 0mm of vi]  {$v$};
    \node (vi1l) [above right=-3.5mm and 0mm of vi1] {$v$};

    \node (vim)  [below right=-3.5mm and 0mm of vi]  {\circled{!}};
    \node (vi1m) [below right=-3.5mm and 0mm of vi1] {\circled{!}};
  
    \begin{pgfonlayer}{background}
      \node[state] (Ti)  [fit=(vi) (vil) (vim)] {};
      \node[state] (Ti1) [fit=(vi1) (vi1l) (vi1m)] {};
    \end{pgfonlayer}  
    
    \node [below right=-1mm and -1mm of Ti]  {$T_i$};
    \node [below right=-1mm and -1mm of Ti1] {$T_{i+1}$};
    
    \node (q) [below=5mm of Ti]  {$q$};
    \node (p) [below=5mm of Ti1] {$p$};
    \node [left=5mm of q]  {\strut Teilrun:};
    \node [right=10mm of q] {\strut $\cdots\qquad f\qquad \cdots$};
  \end{tikzpicture}
\end{center}

\par\medskip\noindent
\textsfbf{Beweis der HA.}~
Damit $v$ überhaupt mit \circled{!} markiert werden kann,
muss $v$ direkt vor der entsprechenden Anwendung des Schrittes~6
Kinder haben.
Diese wurden in irgendeiner vorigen Anwendung von Schritt~2 erzeugt.
Damit dieser Schritt angewendet worden sein kann,
muss gelten:
%
\begin{itemize}
  \item[$(*)$]
    Zwischen $T_i$ und $T_{i+1}$ gibt es einen Zeitpunkt,
    zu dem der Makrozustand von $v$ einen akzeptierenden Zustand $f \in F$
    enthält.
\end{itemize}
%
Wir betrachten nun den Teilrun $T_i \ldots T_{i+1}$ von $s$.
Sei $U$ ein Safra-Baum zwischen $T_i$ und $T_{i+1}$, so dass
Knoten $v$ in $U$ Bedingung $(*)$ erfüllt,
und seien $T_i = S_k$, $U = S_\ell$ und $T_{i+1} = S_m$
für entsprechende $k,\ell,m$ mit $0 \leq k \leq \ell < m$
(die Zeitpunkte des Vorkommens von $T_i,U,T_{i+1}$ auf dem Run $s$;
mit $k \leq \ell$ ist also auch $U = T_i$ erlaubt).

Seien $M,N,P$ die Makrozustände des Knotens $v$ in $T_i,U,T_{i+1}$:
%
\begin{center}
  \begin{tikzpicture}[%
    node distance=25mm,>=Latex,
    initial text="", initial where=below left,
    every state/.style={ellipse,draw=black,thin,fill=black!5,inner sep=.7mm},
    accepting/.style={double distance=1.5pt, double=white},
    every edge/.style={draw=black,thin}
  ]
    \begin{scope}[
      every node/.style={rectangle,rounded corners,draw=black,fill=white,inner sep=1mm}
    ]    
      \node (vi)                {\strut$M$};
      \node (vu)  [right=of vi] {\strut$N$};
      \node (vi1) [right=of vu] {\strut$P$};
    \end{scope}

    \node (vil)  [above right=-3.5mm and 0mm of vi]  {$v$};
    \node (vul)  [above right=-3.5mm and 0mm of vu]  {$v$};
    \node (vi1l) [above right=-3.5mm and 0mm of vi1] {$v$};

    \node (vim)  [below right=-3.5mm and 0mm of vi]  {\circled{!}};
    \node (vum)  [below right=-3.5mm and 0mm of vu]  {\phantom{\circled{!}}};
    \node (vi1m) [below right=-3.5mm and 0mm of vi1] {\circled{!}};
  
    \begin{pgfonlayer}{background}
      \node[state] (Ti)  [fit=(vi) (vil) (vim)] {};
      \node[state] (U)   [fit=(vu) (vul) (vum)] {};
      \node[state] (Ti1) [fit=(vi1) (vi1l) (vi1m)] {};
    \end{pgfonlayer}  
    
    \node [below right=-1mm and -1mm of Ti]  {$T_i$};
    \node [below right=-1mm and -1mm of U]   {$U$};
    \node [below right=-1mm and -1mm of Ti1] {$T_{i+1}$};
  \end{tikzpicture}
\end{center}
%
Mit den eingeführten Bezeichnungen lässt sich die zu zeigende HA nun so formulieren:
%
\begin{itemize}
  \item[$(**)$]
    Für alle $p \in P$ gibt es ein $q \in M$ und einen endlichen Run $q\ldots p$
    von $\Amc$ auf $\alpha[k,m-1]$, der einen akzeptierenden Zustand enthält.
\end{itemize}
%
Um $(**)$ zu beweisen, betrachten wir zunächst den Spezialfall, dass $U$
der einzige Safra-Baum zwischen $T_i$ und $T_{i+1}$ ist, der $(*)$ erfüllt.
Im nächsten Schritt argumentieren wir dann für den allgemeinen Fall.

Im Spezialfall betrachten wir die endlichen Teilruns $T\ldots U$ und $U\ldots T_{i+1}$
von $s$ auf den Teilwörtern $\alpha[k,\ell-1]$ bzw.\ $\alpha[\ell,m-1]$.
Wir schauen uns dazu genauer [1] die Berechnung des Nachfolger-Baums von $U$
und [2] die Berechnung von $T_{i+1}$ aus seinem Vorgängerbaum $X$ (evtl.\ ist $X=U$)
an je einer bestimmten Stelle der Safra-Konstruktion von $\Delta^d$ an:
%
\begin{itemize}
  \item[{[1]}]
    Während der Berechnung des Übergangs $(U,\alpha_\ell,\cdot) \in \Delta^d$
    wird in \emph{Schritt~2} ein Kind $v'$ von $v$ mit Makrozustand $N \cap F$ erzeugt.
    Dieser Knoten $v'$ bleibt in allen Safrabäumen vor $T_{i+1}$ erhalten,
    weil kein Schritt~6 angewendet wird, denn $v$ ist bis $T_{i+1}$ nicht mit \circled{!} markiert.
  \item[{[2]}]
    Während der Berechnung des Übergangs $(X,\alpha_{m-1},T_{i+1}) \in \Delta^d$
    muss die Bedingung in \emph{Schritt~6} erfüllt sein, da $v$ in $T_{i+1}$ markiert ist.
    Folglich haben direkt vorher die Knoten $v$ und $v'$ denselben Makrozustand
    (denn wir sind im Spezialfall; also werden keine weiteren Kinder von $v$ erzeugt).
\end{itemize}
%
Mit diesen Erkenntnissen kann man den Teilrun $T_i \ldots T_{i+1}$ schematisch so veranschaulichen:
%
\begin{center}
  \begin{tikzpicture}[%
    node distance=40mm,>=Latex,
    initial text="", initial where=below left,
    every state/.style={ellipse,draw=black,thin,fill=black!5,inner sep=.7mm},
    accepting/.style={double distance=1.5pt, double=white},
    every edge/.style={draw=black,thin}
  ]
    \begin{scope}[
      every node/.style={rectangle,rounded corners,draw=black,fill=white,inner sep=1mm}
    ]    
      \node (M)                    {\strut$M$};
      \node (N)  [right=of M]      {\strut$N$};
      \node (P)  [right=of N]      {\strut$P$};
      \node (NF) [below= 5mm of N] {\strut$N \cap F$};
      \node (P') [below= 5mm of P] {\strut$\cdots p \cdots$};
    \end{scope}

    \node (Ml)  [above right=-3.5mm and 0mm of M]  {$v$};
    \node (Nl)  [above right=-3.5mm and 0mm of N]  {$v$};
    \node (Pl)  [above right=-3.5mm and 0mm of P]  {$v$};
    \node (NFl) [above right=-3.5mm and 0mm of NF] {$v'$};
    \node (P'l) [above right=-3.5mm and 0mm of P'] {$v'$};

    \node (Mm)  [below right=-3.5mm and 0mm of M]  {\circled{!}};
  
%    \begin{pgfonlayer}{background}
%      \node[state] (Ti)  [fit=(vi) (vil) (vim)] {};
%      \node[state] (U)   [fit=(vu) (vul) (vum)] {};
%      \node[state] (Ti1) [fit=(vi1) (vi1l) (vi1m)] {};
%    \end{pgfonlayer}  
%    
%    \node [below right=-1mm and -1mm of Ti]  {$T_i$};
%    \node [below right=-1mm and -1mm of U]   {$U$};
%    \node [below right=-1mm and -1mm of Ti1] {$T_{i+1}$};

    \node [above=3mm of M] {$T_i$};
    \node [above=3mm of N] {\begin{tabular}{@{}l@{~~}l@{}}[1] & $U$ nach \\ & Schritt 2\end{tabular}};
    \node [above=3mm of P] {\begin{tabular}{@{}l@{~~}l@{}}[2] & $T_{i+1}$ vor \\ & Schritt 6\end{tabular}};
    
    \node [right=10mm of M] {$\stackrel{\text{{\normalsize $\alpha[k,\ell-1]$}}}{=\!=\!=\!=\!=\!=\!\Longrightarrow}$};
    \node [right=9mm of N] {$\stackrel{\text{{\normalsize $\alpha[\ell,m-1]$}}}{=\!=\!=\!=\!=\!=\!\Longrightarrow}$};
    \node [right=6mm of NF] {$\stackrel{\text{{\normalsize $\alpha[\ell,m-1]$}}}{=\!=\!=\!=\!=\!=\!\Longrightarrow}$};

    \path[-] (N) edge (NF)
             (P) edge (P');
  \end{tikzpicture}
\end{center}
%
Da neue Makrozustände in Schritt~3 (knotenweise Potenzmengenkonstruktion) erzeugt werden,
bedeutet die "`untere Zeile"' des Schemas:
%
\begin{quote}
  Für alle $p \in P$ gibt es ein $p' \in N \cap F$ und einen endlichen Run $p'\dots p$
  von \Amc auf $\alpha[\ell,m-1]$.
\end{quote}
%
Diesen Run kann man durch die "`obere Zeile"' ergänzen:
%
\begin{quote}
  Für alle $p' \in N \cap F$ gibt es ein $q \in M$ und einen endlichen Run $q\dots p'$
  von \Amc auf $\alpha[k,\ell-1]$.
\end{quote}
%
Aus diesen beiden Aussagen erhält man wie gewünscht $(**)$.

Es kann natürlich mehrere solche Runs geben kann;
es genügt aber zu wissen, dass es mindestens einen gibt.
Es ist außerdem zu beachten, dass die "`Leserichtung"'
von $(**)$ "`rückwärts"' ist, also "`für alle $p \in P$ existiert ein $q \in M$ \dots"'
und \emph{nicht} umgekehrt.

\par\smallskip
Nun betrachten wir den allgemeinen Fall, dass es mehrere Zwischenzustände
der Art $U$ mit Eigenschaft $(*)$ gibt.
Dann gibt es auch mehrere Situationen [1], und in [2] werden
die Makrozustände \emph{aller} zugehörigen Kinder vereinigt.
Um nun den gesuchten Run $q\ldots p$ auf $\alpha[k,m-1]$ zu erhalten,
muss man nach der ersten "`passenden"' Situation [1]
von der unteren zur oberen Zeile übergehen,
d.\,h.\ wir erhalten $(**)$, indem wir für $p \in P$
in derjenigen Kopie von [1] in die obere Zeile gehen,
in der dasjenige Kind $v'$ von $v$ erzeugt wird, das zu $p \in P$ führt.
%
\begin{center}
  \small
  \begin{tikzpicture}[%
    node distance=30mm,>=Latex,
    initial text="", initial where=below left,
    every state/.style={ellipse,draw=black,thin,fill=black!5,inner sep=.7mm},
    accepting/.style={double distance=1.5pt, double=white},
    every edge/.style={draw=black,thin},
    skip loop/.style={to path={-- ++(0,#1) -| (\tikztotarget)},rounded corners}
  ]
    \begin{scope}[
      every node/.style={rectangle,rounded corners,draw=black,fill=white,inner sep=1mm}
    ]    
      \node[thick] (M)                            {\strut{\boldmath $M$}};
      \node        (N1)     [right=12mm of M]     {\strut$N_1$};
      \node[thick] (N2)     [right=27mm of N1]    {\strut{\boldmath $N_2$}};
      \node        (N3)     [right=24mm of N2]    {\strut$N_3$};
      \node        (P)      [right=46mm of N3]    {\strut$P$};
      \node        (N1F)    [below=5mm of N1]     {\strut$N_1\cap F$};
      \node[thick] (N2F)    [below=5mm of N2]     {\strut{\boldmath $N_2\cap F$}};
      \node        (v1'1')  [left= 7mm of N2F]    {\strut\phantom{$N$}};
      \node        (v2'1'') [below=5mm of N3]     {\strut\phantom{$N$}};
      \node        (v1'1'') [left= 7mm of v2'1''] {\strut\phantom{$N$}};
      \node        (N3F)    [right=7mm of v2'1''] {\strut$N_3\cap F$};
      \node[thick] (v2'2)   [below=5mm of P]      {\strut$\cdots \text{{\boldmath $p$}}\cdots$};
      \node        (v1'2)   [left= 7mm of v2'2]   {\strut\phantom{$N$}};
      \node        (v3'2)   [right=7mm of v2'2]   {\strut\phantom{$N$}};
    \end{scope}

    \node (Ml)      [above right=-3.5mm and 0mm of M] {{\boldmath $v$}};
    \node (N1l)     [right=0mm of N1]                 {$v$};
    \node (N2l)     [right=0mm of N2]                 {{\boldmath $v$}};
    \node (N3l)     [right=0mm of N3]                 {$v$};
    \node (Pl)      [right=0mm of P]                  {$v$};
    \node (N1Fl)    [right=0mm of N1F]                {$v_1'$};
    \node (v1'1'l)  [right=0mm of v1'1']              {$v_1'$};
    \node (N2Fl)    [right=0mm of N2F]                {{\boldmath $v_2'$}};
    \node (v1'1''l) [right=0mm of v1'1'']             {$v_1'$};
    \node (v2'1''l) [right=0mm of v2'1'']             {$v_2'$};
    \node (N3Fl)    [right=0mm of N3F]                {$v_3'$};
    \node (v1'2l)   [right=0mm of v1'2]               {$v_1'$};
    \node (v2'2l)   [right=0mm of v2'2]               {{\boldmath $v_2'$}};
    \node (v3'2l)   [right=0mm of v3'2]               {$v_3'$};

    \node (Mm) [below right=-3.5mm and 0mm of M]  {\circled{!}};
  
%    \begin{pgfonlayer}{background}
%      \node[state] (Ti)  [fit=(vi) (vil) (vim)] {};
%      \node[state] (U)   [fit=(vu) (vul) (vum)] {};
%      \node[state] (Ti1) [fit=(vi1) (vi1l) (vi1m)] {};
%    \end{pgfonlayer}  
%    
%    \node [below right=-1mm and -1mm of Ti]  {$T_i$};
%    \node [below right=-1mm and -1mm of U]   {$U$};
%    \node [below right=-1mm and -1mm of Ti1] {$T_{i+1}$};

    \node [above=13mm of M]  {\normalsize $T_i$};
    \node [above=13mm of N1] {\normalsize $[1]$};
    \node [above=13mm of N2] {\normalsize $[1']$};
    \node [above=13mm of N3] {\normalsize $[1'']$};
    \node [above=13mm of P]  {\normalsize $[2]$};
    
%    \node [right=10mm of M] {$\stackrel{\text{{\normalsize $\alpha[k,\ell-1]$}}}{=\!=\!=\!=\!=\!=\!\Longrightarrow}$};
%    \node [right=9mm of N] {$\stackrel{\text{{\normalsize $\alpha[\ell,m-1]$}}}{=\!=\!=\!=\!=\!=\!\Longrightarrow}$};
%    \node [right=6mm of NF] {$\stackrel{\text{{\normalsize $\alpha[\ell,m-1]$}}}{=\!=\!=\!=\!=\!=\!\Longrightarrow}$};

    \path[-] (N1) edge (N1F)
             (N2) edge (v1'1')
             (N2) edge[thick] (N2F)
             (N3) edge (v1'1'')
             (N3) edge (v2'1'')
             (N3) edge (N3F)
             (P)  edge (v1'2)
             (P)  edge (v2'2)
             (P)  edge (v3'2);

     \path[-,skip loop=-7mm,densely dashed,very thick]
             (N2F) edge[thick] (v2'2);
             
     \path[-,skip loop=7mm,densely dashed,very thick]
             (M) edge[thick] (N2);

     \node [below right=5mm and 1mm of v2'1''] {$\stackrel{\text{{\small $\alpha[\ell',m-1]$}}}{=\!=\!=\!=\!=\!=\!\Longrightarrow}$};
     
     \node [above right=5mm and -3mm of N1] {$\stackrel{\text{{\small $\alpha[k,\ell'-1]$}}}{=\!=\!=\!=\!=\!=\!\Longrightarrow}$};
  \end{tikzpicture}
\end{center}
%
\qedhere

% ===================================================================
\section*{T3.17~ Vollständigkeitsbeweis der Safra-Konstruktion, Details}

\textsfbf{Hilfsaussage [HA].}~
Es gibt einen Knotennamen $v$, für den gilt:
\begin{itemize}
  \item[(a)]
    $\exists m \geqslant 0$ : 
    $S_i$ enthält Knoten $v$ für alle $i \geqslant m$
  \item[(b)]
    $v$ ist in $\infty$ vielen $S_i$ mit $\circled{!}$ markiert
\end{itemize}

\par\medskip\noindent
(Diese Aussage entspricht genau der Akzeptanzbedingung $\Pmc^d$.)

\par\medskip\noindent
\textsfbf{Skizze:}~
$S_0,S_1,\dots,\underbrace{S_m,S_{m+1},S_{m+2},\dots}_{%
  \text{%
    \begin{tabular}{@{}l@{}}
      $v$ in allen $S_i$ enthalten \\
      und unendlich oft markiert
    \end{tabular}
  }
}$

\par\medskip\noindent
\textsfbf{Beweis der HA.}~
Nach Konstruktion enthält der Makrozustand, der im Safrabaum $S_i$ zu Knoten~1 gehört,
alle Zustände von \Amc, die von einem Anfangszustand $q \in I$ aus erreicht werden können,
indem die ersten $i$ Zeichen von $\alpha$ gelesen werden (Potenzmengen\-konstruktion).
Deshalb enthält Knoten~1 in $S_i$ immer den Zustand $q_i$ aus dem Run $r$;
somit hat Knoten~1 immer einen nichtleeren Makrozustand
und wird nie in Schritt~5 entfernt.
Damit erfüllt Knoten~1 die Bedingung~(a) aus der Hilfsaussage.
Wenn er auch Bedingung~(b) erfüllt, ist der Beweis erbracht.

Anderenfalls gibt es einen Zeitpunkt $m' \geq 0$, % (o.\,B.\,d.\,A.\ $m' \geq m$),
so dass Knoten~1 in allen $S_i$ \emph{nicht} mit \circled{!} markiert ist:
%
\begin{center}
  $S_0,S_1,\dots,S_m,S_{m+1},\dots,\!\!\!\!\!\!\underbrace{S_m',S_{m'+1},\dots,S_p,\dots}_{%
    \text{%
      \begin{tabular}{@{}l@{}}
        Knoten 1 ist in keinem $S_i$ markiert % \\
%        und unendlich oft markiert
      \end{tabular}
    }
  }$
\end{center}
%
Da der Run $r$ erfolgreich ist, gibt es einen Zustand $f \in \textsf{Inf}(r) \cap F$.
Sei $p$ der erste Zeitpunkt des Auftretens von $f$ in $r$ hinter $m'$
(d.\,h.\ $q_p = f$ und $q_i \neq f$ für alle $i$ mit $m < i < p$).
Da $q_p = f$, tritt $f$ im Makrozustand des Knotens~1 in $S_p$ auf.
Folglich wird in Schritt~2 der Berechnung von $S_{p+1}$ ein neues jüngstes Kind
zu Knoten~1 hinzugefügt, dessen Makrozustand $f$ enthält.
Da Knoten~1 für den Rest des Runs unmarkiert bleibt,
wird $q_i$ aus dem Run $r$ für alle $i \geq p+1$ in einem \emph{Kind} von 1 auftreten.
Nach endlich vielen Anwendungen von Schritt~4 muss $q_i$ dauerhaft
in einem \emph{festen} Kind $c$ von 1 bleiben.
Dieses Kind erfüllt also Bedingung~(a) der HA.

Nun kann man das bisherige Argument von 1 auf $c$ übertragen:
entweder erfüllt $c$ auch Bedingung~(b),
oder es gibt ein Kind $c'$, das~(a) erfüllt.
Diese Iteration kann man nicht beliebig oft fortsetzen,
weil die Tiefe eines Safrabaums durch $|Q|$ beschränkt ist.
Folglich muss es einen Nachfahren von 1 geben, der Bedingungen~(a)
und~(b) erfüllt.
\qedhere

% ===================================================================
\section*{T3.18~ NBA für eine Kripke-Struktur}

\begin{center}
  \begin{tabular}{@{}l@{\hspace*{20mm}}l@{}}
    Kripke-Struktur $\Smc$ &
    zugehöriger NBA $\Amc_\Smc$ \\[\baselineskip]
    %
    \begin{tikzpicture}[%
      node distance=25mm,>=Latex,
      initial text="", initial where=above,
      every state/.style={draw=black,thin,fill=black!5,inner sep=1mm,minimum size=6mm},
      accepting/.style={double distance=1.5pt, double=white},
      every edge/.style={draw=black,thin}
    ]
      \node[state,initial] (s0)                                  {$s_0$};
      \node[state,initial] (s1) [right of=s0]                    {$s_1$};
      \node[state]         (s2) [below right=10mm and 9mm of s0] {$s_2$};
      
      \node[left =0mm of s0] {$\{p,q\}$};
      \node[right=0mm of s1] {$\{p\}$};
      \node[right=0mm of s2] {$\{q\}$};
      
%      \node[above left=2mm and 12mm of s0] {$\Smc$};
      
      \path[->] (s0) edge [bend left=15] (s1)
                (s1) edge [bend left=15] (s0)
                (s1) edge                (s2)
                (s2) edge                (s0);
    \end{tikzpicture}
    &
    \begin{tikzpicture}[%
      node distance=32mm,>=Latex,
      initial text="", initial where=above,
      every state/.style={draw=black,thin,fill=black!5,inner sep=1mm,minimum size=6mm},
      accepting/.style={double distance=1.5pt, double=white},
      every edge/.style={draw=black,thin}
    ]
      \node[state,initial,accepting] (q0)                                  {$q_0$};
      \node[state,accepting]         (s0) [below left =10mm and 9mm of q0] {$s_0$};
      \node[state,accepting]         (s1) [right of=s0]                    {$s_1$};
      \node[state,accepting]         (s2) [below right=10mm and 9mm of s0] {$s_2$};
      
%      \node[above left=2mm and 12mm of s0] {$\Amc_\Smc$};
      
      \path[->] (q0) edge                node [pos=.2,left =1mm]    {$\{p,q\}$} (s0)
                (q0) edge                node [pos=.2,right=1mm]    {$\{p\}$}   (s1)
                (s0) edge [bend left=10] node [above=-.2mm]         {$\{p\}$}   (s1)
                (s1) edge [bend left=10] node [below=-.2mm]         {$\{p,q\}$} (s0)
                (s1) edge                node [pos=.75,right=1.5mm] {$\{q\}$}   (s2)
                (s2) edge                node [pos=.2,left=1mm]     {$\{p,q\}$} (s0);
    \end{tikzpicture}
  \end{tabular}
\end{center}

% ===================================================================
\section*{T3.19~ NBAs für Beispiel-Eigenschaften}

\subsection*{(1) Mikrowellen-Beispiel}

Hier ist das Alphabet $\Sigma = 2^{\{\textsf{S},\textsf{C},\textsf{H},\textsf{E}\}}$.

\begin{itemize}
  \item[(a)]
    "`Wenn ein Fehler auftritt, dann ist er nach endlicher Zeit behoben."'
    
    Schematisch muss der Automat so aussehen:
    %
    \begin{center}
      \begin{tikzpicture}[%
        node distance=45mm,>=Latex,
        initial text="", initial where=below left,
        every state/.style={draw=black,thin,fill=black!5,inner sep=1mm,minimum size=6mm},
        accepting/.style={double distance=1.5pt, double=white},
        every edge/.style={draw=black,thin}
      ]
        \node[state,initial,accepting] (q0) {};
        \node[state]                   (q1) [right of=q0] {};
        
        \path[->] (q0) edge [bend left=10] node [above] {"`$E$"'}       (q1)
                  (q1) edge [bend left=10] node [below] {"`kein $E$"'}  (q0)
                  (q0) edge [loop below]   node [below] {"`kein $E$"'}  ()
                  (q1) edge [loop below]   node [below] {"`irgendwas"'} ();
      \end{tikzpicture}
    \end{center}
    %
    Dabei steht z.\,B.\ die Beschriftung "`$E$"' für alle Alphabetzeichen
    (Teilmengen von $\{\textsf{S},\textsf{C},\textsf{H},\textsf{E}\}$), die $E$ enthalten,
    und "`kein $E$"' für alle Alphabetzeichen, die $E$ nicht enthalten.
    Sei also $M = \{\{E\}, \{S,E\}, \{C,E\}, \dots, \{S,C,H,E\}\}$
    und $M' = \Sigma \setminus M$. Dann sind die korrekten Kantenbeschriftungen
    im Automaten wie folgt.
    %
    \begin{center}
      \begin{tikzpicture}[%
        node distance=45mm,>=Latex,
        initial text="", initial where=below left,
        every state/.style={draw=black,thin,fill=black!5,inner sep=1mm,minimum size=6mm},
        accepting/.style={double distance=1.5pt, double=white},
        every edge/.style={draw=black,thin}
      ]
        \node[state,initial,accepting] (q0) {};
        \node[state]                   (q1) [right of=q0] {};
        
        \path[->] (q0) edge [bend left=10] node [above] {$M$}      (q1)
                  (q1) edge [bend left=10] node [below] {$M'$}     (q0)
                  (q0) edge [loop below]   node [below] {$M'$}     ()
                  (q1) edge [loop below]   node [below] {$\Sigma$} ();
      \end{tikzpicture}
    \end{center}
  \item[(b)]
    "`Wenn die Mikrowelle gestartet wird,
    fängt sie nach endlicher Zeit an zu heizen."'
    
    Hier nur die schematische Repräsentation des Automaten;
    die korrekten Kanten\-beschriftungen erhält man wie in~(a).
    %
    \begin{center}
      \begin{tikzpicture}[%
        node distance=45mm,>=Latex,
        initial text="", initial where=below left,
        every state/.style={draw=black,thin,fill=black!5,inner sep=1mm,minimum size=6mm},
        accepting/.style={double distance=1.5pt, double=white},
        every edge/.style={draw=black,thin}
      ]
        \node[state,initial,accepting] (q0) {};
        \node[state]                   (q1) [right of=q0] {};
        
        \path[->] (q0) edge [bend left=10] node [above] {"`$S$"'}       (q1)
                  (q1) edge [bend left=10] node [below] {"`$H$"'}       (q0)
                  (q0) edge [loop below]   node [below] {"`kein $S$"'}  ()
                  (q1) edge [loop below]   node [below] {"`irgendwas"'} ();
      \end{tikzpicture}
    \end{center}
  \item[(c)]
    "`Wenn die Mikrowelle gestartet wird,
    ist es \emph{möglich}, danach zu heizen."'
    
    Der Automat ist derselbe wie in~(b),
    nur muss man hier existenzielles Model-Checking
    statt universellem verwenden.
    
\end{itemize}

\subsection*{(2) Nebenläufige Programme}

Hier ist das Alphabet $\Sigma = 2^{\{0,1,10,11,\dots,23\}}$.
%
\begin{itemize}
  \item[(d)]
    "`Es kommt nie vor,
    dass beide Teilprogramme zugleich im kritischen Bereich sind."'

    Schematisch muss der Automat so aussehen:
    %
    \begin{center}
      \begin{tikzpicture}[%
        node distance=45mm,>=Latex,
        initial text="", initial where=below left,
        every state/.style={draw=black,thin,fill=black!5,inner sep=1mm,minimum size=6mm},
        accepting/.style={double distance=1.5pt, double=white},
        every edge/.style={draw=black,thin}
      ]
        \node[state,initial,accepting] (q0) {};
        
        \path[->] (q0) edge [loop below]   node [below] {"`nicht 12 oder nicht 22"'}  ();
      \end{tikzpicture}
    \end{center}
    %
    Dabei steht die Beschriftung "`nicht 12 oder nicht 22"' der Schleife für alle Alphabet\-zeichen
    (Teilmengen von $\{0,1,10,11,\dots,23\}$), die nicht gleichzeitig 12 und 22 enthalten.
    Damit ist die korrekte Kantenbeschriftung für die Schleife
    die Menge $\{X \subseteq \{0,1,10,11,\dots,23\} \mid \{12,22\} \nsubseteq X\}$.
    %
  \item[(e)]
    "`Jedes Teilprogramm kommt beliebig oft in seinen kritischen Bereich."'
    
    Der Automat $\Amc_1$, der beschreibt, dass $P_1$ beliebig oft in seinen kritischen Bereich kommt,
    ist dem in (a) sehr ähnlich (hier wieder nur die schematische Darstellung):
    %
    \begin{center}
      \begin{tikzpicture}[%
        node distance=45mm,>=Latex,
        initial text="", initial where=below left,
        every state/.style={draw=black,thin,fill=black!5,inner sep=1mm,minimum size=6mm},
        accepting/.style={double distance=1.5pt, double=white},
        every edge/.style={draw=black,thin}
      ]
        \node[state,initial,accepting] (q0) {};
        \node[state]                   (q1) [right of=q0] {};
        
        \path[->] (q0) edge [bend left=10] node [above] {"`nicht 12"'}  (q1)
                  (q1) edge [bend left=10] node [below] {"`12"'}        (q0)
                  (q0) edge [loop below]   node [below] {"`12"'}        ()
                  (q1) edge [loop below]   node [below] {"`irgendwas"'} ();
      \end{tikzpicture}
    \end{center}
    %
    Der Automat $\Amc_2$ für $P_2$ ist analog.
    Um zu beschreiben, dass \emph{beide} Programme beliebig oft in den
    jeweiligen kritischen Bereich kommen, muss man den Produktautomaten von $\Amc_1$ und $\Amc_2$ bilden.
    
  \item[(e)]
    "`Jedes Teilprogramm \emph{kann} beliebig oft in seinen kritischen Bereich gelangen."'
    
    Wie (d), aber mit existenziellem Model-Checking.
\end{itemize}

\goodbreak
% ===================================================================
\section*{T3.20~ Beispiele für LTL-Syntax und Semantik}

Betrachte folgenden Pfad $\pi$.
%
\begin{center}
  \begin{tikzpicture}[%
    node distance=15mm,>=Latex,
    initial text="", initial where=below left,
    every state/.style={draw=black,thin,fill=black!5,inner sep=1mm,minimum size=6mm},
    accepting/.style={double distance=1.5pt, double=white},
    every edge/.style={draw=black,thin}
  ]
    \node[state] (q0)               {0};
    \node[state] (q1) [right of=q0] {1};
    \node[state] (q2) [right of=q1] {2};
    \node[state] (q3) [right of=q2] {3};
    \node[state] (q4) [right of=q3] {4};
    \node[state] (q5) [right of=q4] {5};
    \node[state] (q6) [right of=q5] {6};
    \node[state] (q7) [right of=q6] {7};
    \node        (q8) [right of=q7] {$\cdots$};
    
    \node         [below=-.5mm of q1]    {\strut $a$};
    \node         [below=-.5mm of q2]    {\strut $a$};
    \node         [below=-.5mm of q3]    {\strut $a$};
    \node         [below=-.5mm of q4]    {\strut $b$};
    \node         [below=-.5mm of q6]    {\strut $b$};
    \node         [below=-.5mm of q7]    {\strut $a$};
    \node (lastb) [below=0mm of q8]    {\strut $b$};
    \node         [right=2mm of lastb] {\strut $a$~~~$\cdots$};
    
    \path[->] (q0) edge (q1)
              (q1) edge (q2)
              (q2) edge (q3)
              (q3) edge (q4)
              (q4) edge (q5)
              (q5) edge (q6)
              (q6) edge (q7)
              (q7) edge (q8);
  \end{tikzpicture}
\end{center}
%
Das heißt also, $\pi(0) = \emptyset$, $\pi(1) = \{a\}$ usw.
Die Beschriftung "`$b~a~\cdots$"' am rechten Rand bedeutet,
dass für alle $i \geq 4$ gilt: $\pi(2i) = \{b\}$ und $\pi(2i+1) = \{a\}$.
Dann gilt:
%
\begin{xalignat*}{4}
  \pi,0 & \not\models a         & \pi,0 & \models \F a             & \pi,0 & \not\models \G(a \lor b) & \pi,1 & \models a \U b \\
  \pi,0 & \models \lnot a       & \pi,0 & \models \F b             & \pi,6 & \models \G(a \lor b)     & \pi,0 & \not\models a \U b \\
  \pi,0 & \models \X a          & \pi,4 & \models \F a             & \pi,5 & \not\models \G(a \lor b) & \\
  \pi,0 & \not\models \X\lnot a & \pi,0 & \models \X(a \land \F b) & \pi,0 & \models \G\F a           &
\end{xalignat*}
%
Dabei ist $\G\F\varphi$ laut Semantik gleichbedeutend mit "`unendlich oft $\varphi$"'.

\par\medskip
Betrachte nun den folgenden Pfad $\pi'$.
%
\begin{center}
  \begin{tikzpicture}[%
    node distance=15mm,>=Latex,
    initial text="", initial where=below left,
    every state/.style={draw=black,thin,fill=black!5,inner sep=1mm,minimum size=6mm},
    accepting/.style={double distance=1.5pt, double=white},
    every edge/.style={draw=black,thin}
  ]
    \node[state] (q0)               {0};
    \node[state] (q1) [right of=q0] {1};
    \node[state] (q2) [right of=q1] {2};
    \node[state] (q3) [right of=q2] {3};
    \node[state] (q4) [right of=q3] {4};
    \node[state] (q5) [right of=q4] {5};
    \node[state] (q6) [right of=q5] {6};
    \node[state] (q7) [right of=q6] {7};
    \node        (q8) [right of=q7] {$\cdots$};
    
    \node [below=-.5mm of q1] {\strut $a$};
    \node [below=-.5mm of q3] {\strut $a$};
    \node [below=-.5mm of q5] {\strut $a$};
    \node [below=-.5mm of q6] {\strut $b$};
    
    \path[->] (q0) edge (q1)
              (q1) edge (q2)
              (q2) edge (q3)
              (q3) edge (q4)
              (q4) edge (q5)
              (q5) edge (q6)
              (q6) edge (q7)
              (q7) edge (q8);
  \end{tikzpicture}
\end{center}
%
Dann gilt $\pi',0 \not\models a \U b$, aber $\pi',0 \models (\X a \lor \X\X a) \U (\X b)$.

% ===================================================================
\section*{T3.21~ Beispiele für die Erweiterung von Pfaden}

Sei z.\,B. $\varphi_E = \X(a\U b)$ und der Pfad $\pi=s_0s_1s_2\cdots$
wie folgt gegeben:
%
\begin{center}
  \begin{tikzpicture}[%
    node distance=18mm,>=Latex,
    initial text="", initial where=below left,
    every state/.style={draw=black,thin,fill=black!5,inner sep=1mm,minimum size=6mm},
    accepting/.style={double distance=1.5pt, double=white},
    every edge/.style={draw=black,thin}
  ]
    \node[state] (q0)               {0};
    \node[state] (q1) [right of=q0] {1};
    \node[state] (q2) [right of=q1] {2};
    \node[state] (q3) [right of=q2] {3};
    \node[state] (q4) [right of=q3] {4};
    \node[state] (q5) [right of=q4] {5};
    \node[state] (q6) [right of=q5] {6};
    \node        (q7) [right of=q6] {$\cdots$};
    
    \node [below=-.5mm of q0] {\strut $a$};
    \node [below=-.5mm of q1] {\strut $a$};
    \node [below=-.5mm of q2] {\strut $b$};
    \node [below=-.5mm of q4] {\strut $a$};
    \node [below=-.5mm of q5] {\strut $a$};
    \node [below=-.5mm of q6] {\strut $b$};
    
    \path[->] (q0) edge (q1)
              (q1) edge (q2)
              (q2) edge (q3)
              (q3) edge (q4)
              (q4) edge (q5)
              (q5) edge (q6)
              (q6) edge (q7);
  \end{tikzpicture}
\end{center}
%
Das heißt also $s_0 = s_1 = \{a\}$, $s_2 = \{b\}$, $s_3 = \emptyset$ usw.;
dieser Pfad entspricht somit dem Eingabewort, das aus den Zeichen
$\{a\},\{b\},\emptyset,\{a\},\{b\},\emptyset,\dots$ besteht
(jede Teilmenge von \textsf{AV} ist ein Zeichen!).

Der zugehörige erweiterte Pfad $\overline{\pi}=t_0t_1t_2\cdots$ ist der folgende.
%
\begin{center}
  \begin{tikzpicture}[%
    node distance=18mm,>=Latex,
    initial text="", initial where=below left,
    every state/.style={draw=black,thin,fill=black!5,inner sep=1mm,minimum size=6mm},
    accepting/.style={double distance=1.5pt, double=white},
    every edge/.style={draw=black,thin}
  ]
    \node[state] (q0)               {0};
    \node[state] (q1) [right of=q0] {1};
    \node[state] (q2) [right of=q1] {2};
    \node[state] (q3) [right of=q2] {3};
    \node[state] (q4) [right of=q3] {4};
    \node[state] (q5) [right of=q4] {5};
    \node[state] (q6) [right of=q5] {6};
    \node        (q7) [right of=q6] {$\cdots$};
    
    \node [below=-.5mm of q0] {\begin{tabular}{@{}ccc@{}}\strut $a$\\$a\U b$\\$\X(a\U b)$\\$\lnot b$\end{tabular}};
    \node [below=-.5mm of q1] {\begin{tabular}{@{}ccc@{}}\strut $a$\\$a\U b$\\$\X(a\U b)$\\$\lnot b$\\$\vdots$\end{tabular}};
    \node [below=-.5mm of q2] {\begin{tabular}{@{}ccc@{}}\strut $b$\\$a\U b$\\$\lnot\X(a\U b)$\\$\lnot a$\\$\vdots$\end{tabular}};
    \node [below=-.5mm of q3] {\begin{tabular}{@{}ccc@{}}\strut~\\$\lnot(a\U b)$\\$\X(a\U b)$\\$\lnot a,\lnot b$\\$\vdots$\end{tabular}};
    \node [below=-.5mm of q4] {\begin{tabular}{@{}ccc@{}}\strut $a$\\$a\U b$\\$\X(a\U b)$\\$\lnot b$\\$\vdots$\end{tabular}};
    \node [below=-.5mm of q5] {\begin{tabular}{@{}ccc@{}}\strut $a$\\$a\U b$\\$\X(a\U b)$\\$\lnot b$\\$\vdots$\end{tabular}};
    \node [below=-.5mm of q6] {\begin{tabular}{@{}ccc@{}}\strut $b$\\$a\U b$\\$\lnot\X(a\U b)$\\$\lnot a$\\$\vdots$\end{tabular}};
    
    \path[->] (q0) edge (q1)
              (q1) edge (q2)
              (q2) edge (q3)
              (q3) edge (q4)
              (q4) edge (q5)
              (q5) edge (q6)
              (q6) edge (q7);
  \end{tikzpicture}
\end{center}
%
Das heißt also, dass z.\,B.\ $t_0$ die elementare Formelmenge $\{a,\,a\U b,\,\X(a\U b),\,\lnot b\} \subseteq \textsf{cl}(\varphi_E)$ ist.


% ===================================================================
\section*{T3.22~ Beispiele für elementare Formelmengen}

Sei $\varphi_E = a\U(\lnot a \land b)$.
Dann ist $\textsf{cl}(\varphi_E) = \{a,\, \lnot a,\, b,\, \lnot b,\, \lnot a \land b,\, \lnot(\lnot a \land b),\, \varphi_E,\, \lnot \varphi_E\}$.
%
\begin{itemize}
  \item
    $\{a,b,\varphi_E\}$ ist konsistent bezüglich der Aussagenlogik und lokal konsistent bezüglich des \U-Operators,
    aber nicht maximal, weil weder $\lnot a \land b$ noch $\lnot(\lnot a \land b)$ enthalten ist.
  \item
    Fügt man $\lnot a \land b$ hinzu, so ist die resultierende Menge $\{a,\,b,\,\lnot a \land b,\,\varphi_E\}$ 
    zwar maximal, aber nicht mehr konsistent bezüglich der Aussagenlogik,
    da wegen $\lnot a \land b$ auch $\lnot a$ enthalten sein müsste.
  \item
    Fügt man stattdessen $\lnot(\lnot a \land b)$ hinzu, 
    so ist die resultierende Menge \mbox{$\{a,\,b,\,\lnot(\lnot a \land b),\,\varphi_E\}$}
    maximal und konsistent bezüglich der Aussagenlogik,
    aber nicht mehr konsistent bezüglich \U,
    weil nun zwar $\varphi_E = a \U (\lnot a \land b)$ enthalten
    aber weder $\lnot a \land b$ noch $a$ enthalten sind.
  \item
    Die elementaren Formelmengen sind folgende.
    %
    \[
      \begin{array}{@{\{}r@{,~}r@{,~}r@{,~}r@{\}}}
        a       & b       & \lnot(\lnot a \land b) & \varphi_E      \\
        a       & b       & \lnot(\lnot a \land b) & \lnot\varphi_E \\
        a       & \lnot b & \lnot(\lnot a \land b) & \varphi_E      \\
        a       & \lnot b & \lnot(\lnot a \land b) & \lnot\varphi_E \\
        \lnot a & \lnot b & \lnot(\lnot a \land b) & \lnot\varphi_E \\
        \lnot a & b       & \lnot a \land b        & \varphi_E
      \end{array}
    \]
\end{itemize}

% ===================================================================
\section*{T3.23~ Skizzen zur Def.\ der Überführungsrelation des GNBA}

Bedingung \circled{1} besagt, dass im Beispiel in T3.21 \emph{höchstens}
die Transitionen $(t_0,\{a\},t_1)$, $(t_1,\{a\},t_2)$, $(t_2,\{b\},t_3)$, $(t_3,\emptyset,t_4)$ usw.\
erlaubt sind (sofern jeweils Bedingungen \circled{2} und \circled{3} auch erfüllt sind).

\par\medskip
Bedingung \circled{2} kann man so veranschaulichen:
%
\begin{center}
  \begin{tikzpicture}[%
    node distance=27mm,>=Latex,
    initial text="", initial where=below left,
    every state/.style={draw=black,thin,fill=black!5,inner sep=1mm,minimum size=6mm},
    accepting/.style={double distance=1.5pt, double=white},
    every edge/.style={draw=black,thin}
  ]
    \node[state] (q0)               {0};
    \node        (q1)         [right of=q0] {$\cdots$};
    \node[state] (q2)         [right of=q1] {$i$};
    \node[state,ellipse] (q3) [right of=q2] {$i+1$};
    \node        (q4)         [right of=q3] {$\cdots$};
    
    \node (ilabel)  [below=0 of q2] {\begin{tabular}{@{}ccc@{}}\strut $t_i$ \\ enthält \\ $\X\psi$\end{tabular}};
    \node (i1label) [below=0 of q3] {\begin{tabular}{@{}ccc@{}}\strut $t_{i+1}$ \\ enthält \\ $\psi$\end{tabular}};
    \node           [right=-1mm of ilabel] {\emph{gdw.}};
    
    \path[->] (q0) edge (q1)
              (q1) edge (q2)
              (q2) edge (q3)
              (q3) edge (q4);
  \end{tikzpicture}
\end{center}

\par\medskip
Bedingung \circled{3} kann man so veranschaulichen:
%
\begin{center}
  \begin{tikzpicture}[%
    node distance=27mm,>=Latex,
    initial text="", initial where=below left,
    every state/.style={draw=black,thin,fill=black!5,inner sep=1mm,minimum size=6mm},
    accepting/.style={double distance=1.5pt, double=white},
    every edge/.style={draw=black,thin}
  ]
    \node[state] (q0)               {0};
    \node        (q1)         [right of=q0] {$\cdots$};
    \node[state] (q2)         [right of=q1] {$i$};
    \node[state,ellipse] (q3) [right of=q2] {$i+1$};
    \node        (q4)         [right of=q3] {$\cdots$};
    
    \node (ilabel)   [below=0 of q2]         {\begin{tabular}{@{}ccc@{}}\strut $t_i$ \\ enthält \\ $\psi_1\U\psi_2$\end{tabular}};
    \node            [right=-1mm of ilabel]  {\emph{gdw.:}};
    \node (i1label)  [below=0 of q3] {\strut $t_{i+1}$};
    
    \node (ilabel1)  [below=3mm of ilabel]   {\begin{tabular}{@{}ccc@{}}\strut enthält \\ $\psi_2$\end{tabular}};
    \node            [right=-1mm of ilabel1] {\emph{oder}};
    \node (ilabel2)  [below=3mm of ilabel1]  {\begin{tabular}{@{}ccc@{}}\strut enthält \\ $\psi_1$\end{tabular}};
    \node            [right=-1mm of ilabel2] {\emph{und}};
    \node (i1label2) [right of=ilabel2]      {\begin{tabular}{@{}ccc@{}}\strut enthält \\ $\psi_1\U\psi_2$\end{tabular}};
    
    \path[->] (q0) edge (q1)
              (q1) edge (q2)
              (q2) edge (q3)
              (q3) edge (q4);
  \end{tikzpicture}
\end{center}
%
Diese Bedingung nutzt die semantische Äquivalenz
$\psi_1 \U \psi_2 \equiv \psi_2 \lor (\psi_1 \land \X(\psi_1 \U \psi_2))$
(weist diese selbst nach :--)).

% ===================================================================
\section*{{\boldmath T3.24~ GNBA für die Beispiel-Formel $\X a$}}

\begin{itemize}
  \item
    $\varphi = \X a$
  \item
    $\textsf{cl}(\varphi) = \{a, \lnot a, \X a, \lnot \X a\}$
  \item
    Zustände (elementare Formelmengen): da keine $\land$- oder $\U$-Teilformel vorhanden ist,
    sind nur Konsistenz bezüglich $\lnot$ und Maximalität relevant.
    Es gibt also folgende vier elementare Formelmengen:
    %
    \begin{align*}
      t_1 & = \{a,\X a\} \\
      t_2 & = \{a,\lnot\X a\} \\
      t_3 & = \{\lnot a,\X a\} \\
      t_4 & = \{\lnot a,\lnot\X a\}
    \end{align*}
  \item
    Übergänge: es genügt, Bedingung \circled{2} für \X zu überprüfen,
    also gibt es Übergange
    %
    \begin{itemize}
      \item
        von $t_1$ mit $\{a\}$ zu $t_1$ und $t_2$
      \item
        von $t_2$ mit $\{a\}$ zu $t_3$ und $t_4$
      \item
        von $t_3$ mit $\emptyset$ zu $t_1$ und $t_2$
      \item
        von $t_4$ mit $\emptyset$ zu $t_3$ und $t_4$
    \end{itemize}
  \item
    Anfangszustände: $t_1,t_3$ (diese enthalten $\varphi$)
  \item
    Akzeptanzkomponente: $\Fmc = \emptyset$,
    denn es gibt keine $\U$-Teilformeln
    (intuitiv: also braucht auch kein unendlich langes "`Aufschieben"'
    verhindert zu werden). Folglich sind \emph{alle} Runs erfolgreich
    (vgl.\ Definition GNBA).
\end{itemize}
%
Graphische Darstellung:
%
\begin{center}
  \begin{tikzpicture}[%
    node distance=40mm,>=Latex,
    initial text="", initial where=below left,
    every state/.style={ellipse,draw=black,thin,fill=black!5,inner sep=1mm,minimum size=6mm},
    accepting/.style={double distance=1.5pt, double=white},
    every edge/.style={draw=black,thin}
  ]
    \node[state,initial] (t1)                    {$\{a,\X a\}$};
    \node[state]         (t2) [right=30mm of t1] {$\{a,\lnot \X a\}$};
    \node[state,initial] (t3) [below=20mm of t1] {$\{\lnot a,\X a\}$};
    \node[state]         (t4) [below=20mm of t2] {$\{\lnot a,\lnot \X a\}$};
    
    \path[->] (t1) edge                 node [above] {$\{a\}$} (t2)
              (t2) edge                 node [right] {$\{a\}$} (t4)
              (t4) edge                 node [above] {$\emptyset$} (t3)
              (t3) edge                 node [right] {$\emptyset$} (t1)
              (t2) edge [bend right=10] node [above left=0mm and -1mm] {$\{a\}$} (t3)
              (t3) edge [bend right=10] node [below right=0mm and -1mm] {$\emptyset$} (t2)
              (t1) edge [loop above]    node [right] {$~\{a\}$} ()
              (t4) edge [loop below]    node [right] {$~\emptyset$} ();
              
     \node [below right=-1mm and 0mm of t1] {$t_1$};
     \node [below right=-1mm and 0mm of t2] {$t_2$};
     \node [below right=-1mm and 0mm of t3] {$t_3$};
     \node [below right=-1mm and 0mm of t4] {$t_4$};
  \end{tikzpicture}
\end{center}
%


% ===================================================================
\section*{{\boldmath T3.25~ GNBA für die Beispiel-Formel $(\lnot a) \U b$}}

\begin{itemize}
  \item
    $\varphi = (\lnot a) \U b$
  \item
    $\textsf{cl}(\varphi) = \{a, \lnot a, b, \lnot b, (\lnot a) \U b, \lnot((\lnot a) \U b)\}$
  \item
    Zustände (elementare Formelmengen):
    Wegen Konsistenz bezgüglich $\lnot$ und Maximalität
    muss von $a,\lnot a$ bzw.\ $b,\lnot b$ bzw.\
    $(\lnot a) \U b, \lnot((\lnot a) \U b)$
    jeweils genau eine Formel in der Menge enthalten sein.
    Damit gibt es höchstens $2^3 = 8$ elementare Formelmengen.
    Davon sind aber drei nicht konsistent
    bzgl. \U:
    %
    \begin{itemize}
      \item
        $\{a,\,b,\,\lnot((\lnot a) \U b)\}$
        und $\{\lnot a,\,b,\,\lnot((\lnot a) \U b)\}$ \\
        (denn wenn $b \in t$, dann muss auch $\lnot((\lnot a) \U b) \in t$ sein);
      \item
        $\{a,\,\lnot b,\,(\lnot a) \U b\}$\quad
        (denn wenn $(\lnot a) \U b \in t$ und $b \notin t$, dann muss $a \in t$ sein).
    \end{itemize}
    %
    Die verbleibenden fünf Formelmengen sind elementar:
    %
    \begin{xalignat*}{2}
      t_1 & = \{a,\,b,\,(\lnot a) \U b\}                     &
      t_4 & = \{\lnot a,\,\lnot b,\,\lnot ((\lnot a) \U b)\} \\
      t_2 & = \{\lnot a,\,b,\,(\lnot a) \U b\}               &
      t_5 & = \{a,\,\lnot b,\,\lnot ((\lnot a) \U b)\} \\
      t_3 & = \{\lnot a,\,\lnot b,\,(\lnot a) \U b\}
    \end{xalignat*}
  \item
    Übergänge: Bedingung \circled{3} für \U muss eingehalten werden, d.\,h.:
    %
    \begin{center}
     $(\lnot a) \U b \in t$ ~~gdw.~~
     $b \in t$ ~oder~ ($\lnot a \in t$ ~und~ $(\lnot a) \U b \in t'$)
     \qquad $(*)$
    \end{center}
    %
    Nun kann man $(*)$ in folgende logisch äquivalente Aussage umformen,
    die nicht mehr "`gdw."' benutzt:
    %
    \begin{center}
      \begin{tabular}{@{}l@{~~}l@{~~}l@{~~}l@{~~}l@{~~}l@{\qquad}l@{}}
             & $\Big((\lnot a)\U b\in t$    & und & $b \in t\Big)$ &     &                               & (a) \\
        oder & $\Big((\lnot a)\U b\in t$    & und & $\lnot a\in t$ & und & $(\lnot a)\U b \in t'\Big)$   & (b) \\
        oder & $\Big((\lnot a)\U b\notin t$ & und & $b \notin t$   & und & $\lnot a \notin t\Big)$       & (c) \\
        oder & $\Big((\lnot a)\U b\notin t$ & und & $b \notin t$   & und & $(\lnot a)\U b\notin t'\Big)$ & (d)
      \end{tabular}
    \end{center}
    %
    Folglich gibt es Übergänge
    %
    \begin{itemize}
      \item[(a)]
        von $t_1$ mit $\{a,b\}$ zu $t_1,t_2,t_3,t_4,t_5$\quad und \\
        von $t_2$ mit $\{b\}$ zu $t_1,t_2,t_3,t_4,t_5$;
      \item[(b)]
        von $t_3$ mit $\emptyset$ zu $t_1,t_2,t_3$;
      \item[(c)]
        von $t_5$ mit $\{a\}$ zu $t_1,t_2,t_3,t_4,t_5$;
      \item[(d)]
        von $t_4$ mit $\emptyset$ zu $t_4,t_5$.
    \end{itemize}
  \item
    Anfangszustände: $t_1,t_2,t_3$ (diese enthalten $\varphi$)
  \item
    Akzeptanzkomponente:
    %
    \begin{align*}
      \Fmc & = \Big\{\textsf{\textit{M}}_{(\lnot a)\U b}\Big\} \\
           & = \Big\{ \{t \mid (\lnot a)\U b \notin t \text{~oder~} b \in t \} \Big\} \\
           & = \Big\{\{t_1,t_2,t_4,t_5\}\Big\}
    \end{align*}
    %
    Es werden also Runs ausgeschlossen, die irgendwann nur noch durch $t_3$ gehen.
    Auf solchen Runs wird genau $(\lnot a)\U b$ unendlich lange "`aufgeschoben"',
    denn $(\lnot a)\U b \in t_3$ und $b \notin t_3$.
\end{itemize}
%
\enlargethispage*{20mm}
Graphische Darstellung:
%
\vspace*{-2\baselineskip}
\begin{center}
  \begin{tikzpicture}[
     >=Latex,
    initial text="", initial where=above,
    every state/.style={draw=black,thin,fill=black!5,inner sep=1mm,minimum size=6mm},
    accepting/.style={double distance=1.5pt, double=white},
    every edge/.style={draw=black,thin}
  ]
    \node[state,initial,at={(162:40mm)}] (t1) {$t_1$};  
    \node[state,initial,at={( 90:40mm)}] (t2) {$t_2$};  
    \node[state,initial,at={( 18:40mm)}] (t3) {$t_3$};  
    \node[state,at={(306:40mm)}]         (t4) {$t_4$};  
    \node[state,at={(234:40mm)}]         (t5) {$t_5$};
    
    \path[->]
      (t1) edge [loop left]         node [left=0mm]           {$\{a,b\}$}   ()
      (t1) edge [bend left=10]      node [near end,left=2mm]  {$\{a,b\}$}   (t2)
      (t1) edge [bend left=7]       node [pos=.23,above=-1mm] {$\{a,b\}$}   (t3)
      (t1) edge                     node [pos=.45,right=1mm]  {$\{a,b\}$}   (t4)
      (t1) edge [bend left=10]      node [pos=.35,right=-1mm] {$\{a,b\}$}   (t5)
      
      (t2) edge [bend left=10]      node [pos=.50,right=1mm]  {$\{b\}$}     (t1)
      (t2) edge [out=60,in=30,loop] node [right=0mm]          {$\{b\}$}     ()
      (t2) edge [bend left=10]      node [right=1mm]          {$\{b\}$}     (t3)
      (t2) edge                     node [near end,right=0mm] {$\{b\}$}     (t4)
      (t2) edge [bend left=7]       node [pos=.20,right=-1mm]  {$\{b\}$}     (t5)
      
      (t3) edge [bend left=7]       node [pos=.20,below]      {$\emptyset$} (t1)
      (t3) edge [bend left=10]      node [pos=.35,left=2mm]   {$\emptyset$} (t2)
      (t3) edge [loop right]        node [right=0mm]          {$\emptyset$} ()

      (t4) edge [loop right]        node [right=0mm]          {$\emptyset$} ()
      (t4) edge [bend left=10]      node [below]              {$\emptyset$} (t5)

      (t5) edge [bend left=10]      node [pos=.50,left=0mm]   {$\{a\}$}     (t1)
      (t5) edge [bend left=7]       node [pos=.25,left=-1mm]  {$\{a\}$}     (t2)
      (t5) edge                     node [pos=.25,left=0mm]   {$\{a\}$}     (t3)
      (t5) edge [bend left=7]       node [pos=.5,above=-1mm]  {$\{a\}$}     (t4)
      (t5) edge [loop left]         node [left=0mm]           {$\{a\}$}     ()
      ;  
  \end{tikzpicture}
\end{center}

% ===================================================================
% ===================================================================
% ===================================================================
\part[Endliche Automaten auf unendlichen Bäumen]{Endliche Automaten \\ auf unendlichen Bäumen}

% ===================================================================
\section*{T4.1~ LTL-Formeln "`zu stark/schwach"'}

Sei $\varphi_1 := \G(e \to \F \neg e)$ und $\varphi_2 := \G\F\neg e$.

\paragraph*{{\boldmath $\varphi_1,\varphi_2$ "`zu stark für universelles Model Checking"'.}}
%Sei $\varphi_1 := \G(e \to \F \neg e)$ und $\varphi_2 := \G\F\neg e$.
Wir betrachten folgende Kripke-Struktur $\Smc$:
%
\begin{center}
  \begin{tikzpicture}[%
    node distance=20mm,>=Latex,
    initial text="", initial where=below left,
    every state/.style={draw=black,thin,fill=black!5,inner sep=1mm,minimum size=6mm},
    accepting/.style={double distance=1.5pt, double=white},
    every edge/.style={draw=black,thin}
  ]
    \node[state,initial] (s0)                   {$s_0$};
    \node[state]         (s1) [right of=s0]     {$s_1$};
    \node[state]         (s2) [below=8mm of s0] {$s_2$};
    
    \node [above=0mm of s0] {$e$};
%    \node [right=0mm of s2] {$\lnot e$};
    \node [above=0mm of s1] {$e$};
    
    \path[->] (s0) edge[bend left=15] (s1)
              (s1) edge[bend left=15] (s0)
              (s2) edge (s0)
              (s1) edge (s2);
  \end{tikzpicture}
\end{center}
%
Intuitiv gesprochen, erfüllt $\Smc$ die Eigenschaft, die mit $\varphi_1$ bzw.\
$\varphi_2$ ausgedrückt werden soll:
es ist \emph{möglich}, nach Besuchen von $s_0$ bzw.\ $s_1$ den Pfad so fortzusetzen,
dass nach endlich vielen Schritten $s_2$ besucht wird.

Nach der Semantik von LTL wird jedoch ein Pfad "`festgehalten"'.
Wenn wir also z.\,B.\ den Pfad $\pi = (s_0s_1)^\omega$ betrachten,
dann gilt $\pi,0 \not\models \varphi_i$, $i=1,2$
(aber für den Pfad $\pi' = (s_0s_1s_2)^\omega$ gilt
$\pi',0 \models \varphi_i$, $i=1,2$).

\paragraph*{{\boldmath $\varphi_1,\varphi_2$ "`zu schwach für existenzielles Model Checking"'.}}
%Sei wieder $\varphi_1 := \G(e \to \F \neg e)$ und $\varphi_2 := \G\F\neg e$.
Wir betrachten folgende Kripke-Struktur $\Smc'$:
%
\begin{center}
  \begin{tikzpicture}[%
    node distance=20mm,>=Latex,
    initial text="", initial where=below left,
    every state/.style={draw=black,thin,fill=black!5,inner sep=1mm,minimum size=6mm},
    accepting/.style={double distance=1.5pt, double=white},
    every edge/.style={draw=black,thin}
  ]
    \node[state,initial] (s0)                   {$s_0$};
    \node[state]         (s1) [right of=s0]     {$s_1$};
    \node[state]         (s2) [right of=s1]     {$s_2$};
    \node[state]         (s3) [below=8mm of s0] {$s_3$};
    
    \node [above=0mm of s1] {$e$};
%    \node [right=0mm of s2] {$\lnot e$};
    \node [below=0mm of s3] {$e$};
    
    \path[->] (s0) edge (s1)
              (s1) edge (s2)
              (s2) edge [loop right] ()
              (s0) edge (s3)
              (s3) edge [loop right] ();
  \end{tikzpicture}
\end{center}
%
Intuitiv gesprochen, erfüllt $\Smc'$ die gewünschte Eigenschaft
\emph{nicht}: der Pfad $s_0s_3$ kann nicht mehr wie gewünscht fortgesetzt werden.

Nach der Semantik von LTL genügt es jedoch, einen Pfad $\pi$ zu finden,
für den $\pi,0 \models \varphi_i$ gilt,
und das ist z.\,B. $\pi = s_1s_1s_2^\omega$.

\pagebreak
% ===================================================================
\section*{T4.2~ Beispiel-Berechnungsbaum}

Wenn man die Beispielstruktur Mikrowelle im Zustand 3 auf"|faltet,
erhält man:
%
\begin{center}
  \begin{tikzpicture}[%
    node distance=20mm,>=Latex,
    every node/.style={circle,draw=black,thin,fill=black!5,inner sep=1mm,minimum size=6mm},
    level 1/.style = {sibling distance = 35mm, level distance =  6mm},
    level 2/.style = {sibling distance = 20mm, level distance =  9mm},
    level 3/.style = {sibling distance = 10mm, level distance =  9mm},
    level 4/.style = {sibling distance = 10mm, level distance =  9mm},
    edge from parent/.style = {draw=black, thin, -Latex}%
  ]
    \node {3}
    child {
      node {1}
      child {
        node {2}
        child {
          node {5}
          child {
            node (n1) {2}
          }
          child {
            node (n2) {3}
          }
        }
      }
      child {node (n3) {3}}
    }
    child {
      node {6}
      child {
        node {7}
        child {
          node {4}
          child {
            node (n4) {4}
          }
          child {
            node (n5) {3}
          }
          child {
            node (n6) {1}
          }
        }
      }
    }
    ;
    
    \foreach \x in {1,...,6}
      \node [draw=none,fill=none,below=-1mm of n\x] {$\vdots$};
    
  \end{tikzpicture}
\end{center}
%
Die gegebene Struktur $\Smc$ ist eine endliche Repräsentation dieses unendlichen Baums
(und vieler weiterer Bäume).

%\enlargethispage{10mm}
% ===================================================================
\section*{T4.3~ Beweis Teil 2 des Ausdrucksstärke-Lemmas}

\textsfbf{Lemma 4.5~(2).}~
Es gibt keine zu $\F\G p$ äquivalente CTL-Zustandsformel.

\par\medskip\noindent
\textsfbf{Beweis.}~
Betrachte zwei \emph{Folgen} $\Smc_0,\Smc_1,\Smc_2,\dots$ und
$\Smc'_0,\Smc'_1,\Smc'_2,\dots$ von Kripke-Strukturen, die wie folgt aufgebaut sind.
Für $n \geq 0 $ ist
%
\begin{center}
  \begin{tikzpicture}[%
    node distance=10mm,>=Latex,
    initial text="", initial where=below left,
    every state/.style={ellipse,draw=black,thin,fill=black!5,inner sep=1mm,minimum size=6mm},
    accepting/.style={double distance=1.5pt, double=white},
    every edge/.style={draw=black,thin}
  ]
    \node[state,initial]             (sn)                        {\rule{0pt}{7.5pt}$s_n$};
    \node[state]                     (tn)   [right= 6mm of sn]   {\rule{0pt}{7.5pt}$t_n$};
    \node[state]                     (sn-1) [right=12mm of tn]   {\rule{0pt}{7.5pt}$s_{n-1}$};
    \node[state]                     (tn-1) [right= 6mm of sn-1] {\rule{0pt}{7.5pt}$t_{n-1}$};
    \node[state,draw=none,fill=none] (dots) [right=12mm of tn-1] {$\cdots$};
    \node[state]                     (s0)   [right=12mm of dots] {\rule{0pt}{7.5pt}$s_0$};
    \node[state]                     (t0)   [right= 6mm of s0]   {\rule{0pt}{7.5pt}$t_0$};
    
    \node[state,initial]             (sn')   [below=15mm of sn]    {\rule{0pt}{7.5pt}$s_n$};
    \node[state]                     (tn')   [right= 6mm of sn']   {\rule{0pt}{7.5pt}$t_n$};
    \node[state]                     (sn-1') [right=12mm of tn']   {\rule{0pt}{7.5pt}$s_{n-1}$};
    \node[state]                     (tn-1') [right= 6mm of sn-1'] {\rule{0pt}{7.5pt}$t_{n-1}$};
    \node[state,draw=none,fill=none] (dots') [right=12mm of tn-1'] {$\cdots$};
    \node[state]                     (s0')   [right=12mm of dots'] {\rule{0pt}{7.5pt}$s_0$};
    \node[state]                     (t0')   [right= 6mm of s0']   {\rule{0pt}{7.5pt}$t_0$};
    
    \node [above=0mm of tn]   {$p$};
    \node [above=0mm of tn-1] {$p$};
    \node [above=0mm of t0]   {$p$};
    
    \node [above=0mm of tn']   {$p$};
    \node [above=0mm of tn-1'] {$p$};
    \node [above=0mm of t0']   {$p$};
    
    \node [above left=0mm and 5mm of sn]  {$\Smc_n$};
    \node [above left=0mm and 5mm of sn'] {$\Smc'_n$};
    
    \path[->] (sn)   edge [bend left=15] (tn)
              (tn)   edge [bend left=15] (sn)
              (tn)   edge (sn-1)
              (sn-1) edge (tn-1)
              (tn-1) edge (dots)
              (dots) edge (s0)
              (s0)   edge (t0)
              (tn)   edge [loop below] ()
              (tn-1) edge [loop below] ()
              (t0)   edge [loop below] ()
              ;
              
    \path[->] (sn')   edge (tn')
              (tn')   edge (sn-1')
              (sn-1') edge (tn-1')
              (tn-1') edge (dots')
              (dots') edge (s0')
              (s0')   edge (t0')
              (tn')   edge [loop below] ()
              (tn-1') edge [loop below] ()
              (t0')   edge [loop below] ()
              ;
  \end{tikzpicture}
\end{center}
%
Insbesondere unterscheidet sich $\Smc'_n$ von $\Smc_n$ nur durch die fehlende Kante von $t_n$ nach $s_n$.

Nun gilt für alle $n \geq 0$:

\par\medskip\noindent
\begin{tabular}{@{\qquad}lll@{}}
  (i)  & $\Smc_n  \not\models \F\G p$ & \text{(wegen Pfad $(s_nt_n)^\omega$ ab $s_n \in \S_0$)} \\
  (ii) & $\Smc'_n \models \F\G p$     & \text{(weil jeder Pfad auf ein $(t_i)^\omega$ enden muss, für ein $i \leq n$)}
\end{tabular}

\par\medskip\noindent
Außerdem zeigt man leicht per Induktion über $n$ die folgende Aussage:
Für alle $n \geq 0$ und alle CTL-Zustandsformeln $\zeta$ der Länge $\leq n$ gilt:

\par\medskip\noindent
\begin{tabular}{@{\qquad}lll@{}}
  (iii) & $\Smc_n \models \zeta$ ~gdw.~ $\Smc'_n \models \zeta$
\end{tabular}

\par\medskip\noindent
Angenommen, es gebe eine CTL-Zustandsformel $\zeta$ mit $\zeta \equiv \F\G p$.
Sei $n := |\zeta|$.
Dann gilt wegen~(i) und~(ii): $\Smc_n \not\models \zeta$ und $\Smc'_n \models \zeta$.
Das widerspricht aber~(iii).\qedhere

\goodbreak
% ===================================================================
\section*{T4.4~ Beispiel für den Model-Checking-Algorithmus}

\dots\ folgt \dots


%% ===================================================================
%% ===================================================================
%% ===================================================================
%\part*{Anhang}
%\addcontentsline{toc}{part}{Anhang}
%
%% ===================================================================
%\section*{Griechische Buchstaben}
%
%\paragraph*{Kleinbuchstaben}
%~\par%\vspace*{-.2\baselineskip}
%$\alpha$ \dotfill alpha \\
%$\beta$ \dotfill beta \\
%$\gamma$ \dotfill gamma \\
%$\delta$ \dotfill delta \\
%$\epsilon$ \dotfill epsilon \\
%$\zeta$ \dotfill zeta \\
%$\eta$ \dotfill eta \\
%$\vartheta,\theta$ \dotfill theta \\
%$\iota$ \dotfill iota \\
%$\kappa$ \dotfill kappa \\
%$\lambda$ \dotfill lambda\\
%$\mu$ \dotfill my\\
%$\nu$ \dotfill ny \\
%$\xi$ \dotfill xi\\
%$o$ \dotfill omikron\\
%$\pi$ \dotfill pi\\
%$\rho$ \dotfill rho\\
%$\sigma,\varsigma$ \dotfill sigma\\
%$\tau$ \dotfill tau\\
%$\upsilon$ \dotfill ypsilon\\
%$\varphi,\phi$ \dotfill phi\\
%$\chi$ \dotfill chi\\
%$\psi$ \dotfill psi\\
%$\omega$ \dotfill omega
%
%% \newpage
%
%\paragraph*{Großbuchstaben}
%~\par%\vspace*{-.2\baselineskip}
%$\Gamma$ \dotfill Gamma \\
%$\Delta$ \dotfill Delta \\
%$\Theta$ \dotfill Theta \\
%$\Lambda$ \dotfill Lambda \\
%$\Pi$ \dotfill Pi \\
%$\Xi$ \dotfill Xi \\
%$\Sigma$ \dotfill Sigma \\
%$\Upsilon$ \dotfill Ypsilon \\
%$\Phi$ \dotfill Phi \\
%$\Psi$ \dotfill Psi \\
%$\Omega$ \dotfill Omega
%
%Die übrigen griechischen Großbuchstaben werden genauso geschrieben wie die entsprechenden lateinischen.

%\pagebreak
%\addcontentsline{toc}{part}{Literaturverzeichnis}
%\bibliographystyle{babalpha}
%\bibliography{biblio.bib}
\end{document}
